% Chapter 1. All quoted numbers come from chapters/ch01-numbers.tex, written by code/ch01/run_all.py.
% Generated by code/ch01/run_all.py. Do not edit by hand.
\newcommand{\breakEvenYear}{7}
\newcommand{\ccEquivEight}{7.6961}
\newcommand{\chebOne}{100.0}
\newcommand{\chebThree}{11.1}
\newcommand{\chebTwo}{25.0}
\newcommand{\convexPearson}{0.84}
\newcommand{\convexSpearman}{1.00}
\newcommand{\crossRate}{24.03}
\newcommand{\cumTrough}{\ensuremath{-}97.00}
\newcommand{\cumTroughYear}{4}
\newcommand{\divFloorTwo}{8.94}
\newcommand{\divTenTwo}{10.58}
\newcommand{\downLog}{\ensuremath{-}40.55}
\newcommand{\dpiEnd}{1.72}
\newcommand{\dpiFive}{0.39}
\newcommand{\dwrPer}{\ensuremath{-}3.02}
\newcommand{\earA}{8.0000}
\newcommand{\earB}{8.1600}
\newcommand{\earC}{8.2432}
\newcommand{\earD}{8.3000}
\newcommand{\earE}{8.3278}
\newcommand{\earF}{8.3287}
\newcommand{\exBeta}{0.5311}
\newcommand{\exCov}{2.3900}
\newcommand{\exMeanF}{0.6333}
\newcommand{\exMeanM}{0.5000}
\newcommand{\exRho}{0.9917}
\newcommand{\exSdF}{1.1361}
\newcommand{\exSdM}{2.1213}
\newcommand{\exSpear}{1.0000}
\newcommand{\exVarM}{4.5000}
\newcommand{\exkTfive}{6}
\newcommand{\fcA}{22.00}
\newcommand{\fcB}{24.00}
\newcommand{\fcC}{22.00}
\newcommand{\fcD}{17.00}
\newcommand{\fcE}{12.00}
\newcommand{\fcF}{0.00}
\newcommand{\fcG}{0.00}
\newcommand{\fcH}{0.00}
\newcommand{\fcI}{0.00}
\newcommand{\fcJ}{0.00}
\newcommand{\fdA}{0.00}
\newcommand{\fdB}{0.00}
\newcommand{\fdC}{0.00}
\newcommand{\fdD}{0.00}
\newcommand{\fdE}{0.00}
\newcommand{\fdF}{38.23}
\newcommand{\fdG}{42.25}
\newcommand{\fdH}{38.23}
\newcommand{\fdI}{28.17}
\newcommand{\fdJ}{20.11}
\newcommand{\feeDrag}{3.47}
\newcommand{\fnA}{18.00}
\newcommand{\fnB}{39.80}
\newcommand{\fnC}{60.00}
\newcommand{\fnD}{85.92}
\newcommand{\fnE}{112.44}
\newcommand{\fnF}{93.90}
\newcommand{\fnG}{67.79}
\newcommand{\fnH}{41.44}
\newcommand{\fnI}{17.49}
\newcommand{\fnJ}{0.00}
\newcommand{\fundDistrib}{166.99}
\newcommand{\fundIRR}{11.53}
\newcommand{\fundIRRNoFee}{15.00}
\newcommand{\fundPaidIn}{97.00}
\newcommand{\gocBand}{0.030}
\newcommand{\gocBeyondFour}{15}
\newcommand{\gocDN}{4,377}
\newcommand{\gocDW}{1.961}
\newcommand{\gocDexk}{2.35}
\newcommand{\gocDmean}{0.041}
\newcommand{\gocDsd}{4.93}
\newcommand{\gocDskew}{\ensuremath{-}0.08}
\newcommand{\gocDur}{4.77}
\newcommand{\gocEnd}{15 October 2026}
\newcommand{\gocExpectFour}{0.28}
\newcommand{\gocJB}{1,013}
\newcommand{\gocMaxDn}{29}
\newcommand{\gocMaxDnDate}{10 November 2022}
\newcommand{\gocMaxUp}{23}
\newcommand{\gocMaxUpDate}{9 October 2009}
\newcommand{\gocN}{4,378}
\newcommand{\gocRJB}{926}
\newcommand{\gocRatioMonth}{0.964}
\newcommand{\gocRatioQuarter}{0.960}
\newcommand{\gocRatioWeek}{1.025}
\newcommand{\gocRexk}{2.24}
\newcommand{\gocRhoOne}{0.019}
\newcommand{\gocRmean}{1.53}
\newcommand{\gocRsd}{3.68}
\newcommand{\gocRskew}{0.11}
\newcommand{\gocSeMonth}{0.047}
\newcommand{\gocSeQuarter}{0.082}
\newcommand{\gocSha}{40311d546c79}
\newcommand{\gocShaA}{40311d546c79eb75a165898f4dba862b}
\newcommand{\gocShaB}{3f71c6c619b2def680abaa4528aab162}
\newcommand{\gocStart}{1 April 2009}
\newcommand{\gocYmax}{4.42}
\newcommand{\gocYmaxDate}{October 2023}
\newcommand{\gocYmin}{0.30}
\newcommand{\gocYminDate}{August 2020}
\newcommand{\indexCAGR}{5.35}
\newcommand{\interimFive}{9.64}
\newcommand{\interimOne}{\ensuremath{-}28.18}
\newcommand{\interimSeven}{11.53}
\newcommand{\interimThree}{0.67}
\newcommand{\interimTwo}{\ensuremath{-}12.25}
\newcommand{\irB}{8}
\newcommand{\irC}{\ensuremath{-}12}
\newcommand{\irD}{15}
\newcommand{\irE}{10}
\newcommand{\irF}{6}
\newcommand{\irG}{9}
\newcommand{\irH}{\ensuremath{-}4}
\newcommand{\irI}{12}
\newcommand{\irJ}{7}
\newcommand{\irrLong}{25.99}
\newcommand{\irrShort}{30.00}
\newcommand{\irrX}{50.00}
\newcommand{\irrXY}{23.85}
\newcommand{\irrY}{10.00}
\newcommand{\ixA}{100.00}
\newcommand{\ixB}{108.00}
\newcommand{\ixC}{95.04}
\newcommand{\ixD}{109.30}
\newcommand{\ixE}{120.23}
\newcommand{\ixF}{127.44}
\newcommand{\ixG}{138.91}
\newcommand{\ixH}{133.35}
\newcommand{\ixI}{149.35}
\newcommand{\ixJ}{159.81}
\newcommand{\jbCrit}{5.991}
\newcommand{\jbPowFifty}{39.9}
\newcommand{\jbPowFiveHundred}{99.4}
\newcommand{\jbPowHundred}{62.8}
\newcommand{\jbPowSixty}{45.3}
\newcommand{\jbPowThirty}{25.5}
\newcommand{\jbPowThousand}{100.0}
\newcommand{\jbPowTwenty}{16.8}
\newcommand{\jbPowTwoFifty}{91.3}
\newcommand{\jbSizeFifty}{3.8}
\newcommand{\jbSizeFiveHundred}{4.7}
\newcommand{\jbSizeHundred}{4.1}
\newcommand{\jbSizeSixty}{3.9}
\newcommand{\jbSizeThirty}{3.1}
\newcommand{\jbSizeThousand}{4.9}
\newcommand{\jbSizeTwenty}{2.2}
\newcommand{\jbSizeTwoFifty}{4.4}
\newcommand{\jbSnipNorm}{3.24}
\newcommand{\jbSnipNormP}{0.1978}
\newcommand{\jbSnipT}{10.72}
\newcommand{\jbSnipTP}{0.0047}
\newcommand{\ksPME}{1.297}
\newcommand{\lapOutOne}{24.31}
\newcommand{\lapOutThree}{1.44}
\newcommand{\lapOutTwo}{5.91}
\newcommand{\lnExkFive}{3.31}
\newcommand{\lnExkOne}{0.54}
\newcommand{\lnExkTwenty}{32.29}
\newcommand{\lnMeanFive}{1.464}
\newcommand{\lnMeanOne}{1.079}
\newcommand{\lnMeanTwenty}{4.591}
\newcommand{\lnMedianFive}{1.350}
\newcommand{\lnMedianOne}{1.062}
\newcommand{\lnMedianTwenty}{3.320}
\newcommand{\lnModeFive}{1.148}
\newcommand{\lnModeOne}{1.028}
\newcommand{\lnModeTwenty}{1.737}
\newcommand{\lnSkewFive}{1.33}
\newcommand{\lnSkewOne}{0.55}
\newcommand{\lnSkewTwenty}{3.74}
\newcommand{\lnValue}{\ensuremath{-}44.14}
\newcommand{\logFull}{2.372}
\newcommand{\logPart}{8.066}
\newcommand{\logPartLoss}{\ensuremath{-}7.904}
\newcommand{\mirrFVpos}{257.60}
\newcommand{\mirrFundEight}{9.94}
\newcommand{\mirrFundTwelve}{11.04}
\newcommand{\mirrFundTwelveTwelve}{11.74}
\newcommand{\mirrMulti}{9.93}
\newcommand{\mirrPVneg}{213.17}
\newcommand{\multiHi}{20.00}
\newcommand{\multiLo}{10.00}
\newcommand{\noneMinNPV}{10.00}
\newcommand{\normOutOne}{31.73}
\newcommand{\normOutThree}{0.27}
\newcommand{\normOutTwo}{4.55}
\newcommand{\npvLongTen}{50.26}
\newcommand{\npvLongTenAbs}{50.26}
\newcommand{\npvShortTenAbs}{18.18}
\newcommand{\pnlFull}{240,000}
\newcommand{\pnlPart}{210,000}
\newcommand{\polFirst}{0.25}
\newcommand{\polNow}{2.25}
\newcommand{\polStart}{21 April 2009}
\newcommand{\putPrem}{0.394}
\newcommand{\pwExk}{29.01}
\newcommand{\pwExkS}{27.64}
\newcommand{\pwJB}{8546.1}
\newcommand{\pwMean}{0.420}
\newcommand{\pwProbLoss}{8.7}
\newcommand{\pwSD}{0.931}
\newcommand{\pwSkew}{\ensuremath{-}4.96}
\newcommand{\pwSkewS}{\ensuremath{-}4.76}
\newcommand{\rvpiEnd}{0.00}
\newcommand{\rvpiFive}{0.97}
\newcommand{\simpleFull}{2.40}
\newcommand{\simplePart}{8.40}
\newcommand{\simplePartLoss}{\ensuremath{-}7.60}
\newcommand{\smAlpha}{0.6}
\newcommand{\smAnnModel}{13.9}
\newcommand{\smAnnObs}{6.5}
\newcommand{\smAnnTrue}{13.4}
\newcommand{\smBetaFactor}{0.4}
\newcommand{\smCorrFactor}{0.8}
\newcommand{\smDW}{0.855}
\newcommand{\smDWapprox}{0.868}
\newcommand{\smObsSD}{1.88}
\newcommand{\smRatioSample}{0.486}
\newcommand{\smRatioTheory}{0.500}
\newcommand{\smRhoOne}{0.566}
\newcommand{\smTheoryRatio}{50.0}
\newcommand{\smTrueSD}{3.87}
\newcommand{\smUnRho}{\ensuremath{-}0.017}
\newcommand{\smUnSD}{3.87}
\newcommand{\tOutOne}{25.32}
\newcommand{\tOutThree}{1.17}
\newcommand{\tOutTwo}{4.93}
\newcommand{\tailNormFour}{6.3}
\newcommand{\tailTFour}{357}
\newcommand{\tvpiEnd}{1.72}
\newcommand{\tvpiFive}{1.36}
\newcommand{\twrPer}{3.92}
\newcommand{\twrTotal}{8.00}
\newcommand{\udApprox}{\ensuremath{-}0.35}
\newcommand{\udApproxSample}{\ensuremath{-}9.03}
\newcommand{\udHalfVar}{8.68}
\newcommand{\udMean}{8.33}
\newcommand{\upLog}{40.55}
\newcommand{\vrExTen}{1.02}
\newcommand{\vrExThree}{1.07}
\newcommand{\vrExTwo}{1.10}
\newcommand{\vrLimitSix}{4.0}
\newcommand{\vrTwelveSix}{3.376}
\newcommand{\vrTwelveSixSqrt}{1.837}

\begin{chapterrefs}
\chapter{Returns, Rates and the Shape of Risk}
\label{ch:returns}

Every later chapter of this book prices, compares or risk-manages something, and each of those tasks
starts with a number that claims to describe performance. For a listed share that number is easy to
get: two prices and perhaps a dividend. For a private equity fund, a farmland partnership or a futures
programme it is not. Cash goes in and out at irregular dates, some positions have no value when they
are opened, and the values in between come from appraisals. This chapter builds the
arithmetic and the statistics needed to turn such records into returns whose assumptions and limits
are visible, and it shows, case by case, where the familiar tools give the wrong answer.

First we have to establish what was earned. \Crefrange{sec:rates}{sec:notional} get a return
from quoted rates, from cash that moves at irregular dates and from positions that cost nothing to
enter, and ask which of the resulting numbers can be believed. Only then can we ask how much risk a
short record of such returns conceals, which is the subject of \crefrange{sec:distributions}{sec:dependence}. \Cref{sec:case} applies those tools to seventeen years of daily Government of
Canada yields.

% =====================================================================================
\section{What the numbers are for}
\label{sec:purpose}

\begin{keybox}[How to read this chapter]
Most results are proved where they are stated and then checked with a number; a few standard
theorems, such as Descartes' rule of signs, Chebyshev\index{Chebyshev}'s inequality and the central limit theorem, are
used without proof. The worked examples can be done with a calculator. The figures and the
quoted numbers come from the scripts in \texttt{code/ch01/}, and every simulation draws\ix{Simulation} from NumPy's
default generator with seed 101. The scripts share one generator, \texttt{default\_rng(101)}, which
\texttt{run\_all.py} passes from script to script in a fixed order; a script run on its own first
replays the draws made before it, so it prints the same numbers. The short Jarque-Bera\ix{Jarque-Bera test}\index{Jarque-Bera} listing starts
a fresh \texttt{default\_rng(101)}, and the Jarque-Bera run at $n=60$ uses
\texttt{default\_rng([101, 60])}, the same seed with a stream number. A test checks that no script
draws from any other generator.
\end{keybox}

Suppose a pension plan\ix{Pension plan} pays a private equity manager 2\% of its commitment every year plus a fifth
of the profits. Ten years later the trustees want to know what the money bought, and the answer
needs a return, a measure of the risk taken and a fair comparison with what the plan could have
earned elsewhere. A fee that buys nothing, or a risk the trustees did not know they were running,
leaves the plan poorer even when the manager reports a profit. The trustees can only tell which
happened if they have a sound figure for what was earned and for the risk behind it, and producing
such figures is the subject of this chapter.

A farmland partnership is alternative twice over: appraisers set the value of its farms, and the
fund that holds them cannot be redeemed for ten years. There is no single accepted definition of an
alternative investment\ix{Alternative investments}. For this chapter a more useful test is
whether the liquidity, the valuation or the payoff of an investment differs materially from that of
a conventional public portfolio. Many private alternative funds are limited partnerships with negotiated terms, while some
alternatives, such as listed
infrastructure and exchange-traded commodity products, trade on public markets\ix{Public markets}. What they share is
the trouble they make for a performance measure\ix{Performance measurement}: returns that are asymmetric, values that are
appraised instead of traded, and fees that make comparison with a public portfolio difficult.

\FloatBarrier
% =====================================================================================
\section{Compounding and the two kinds of return}
\label{sec:rates}

\subsection{Quoted rates and growth}

Let $P$ be an amount invested for $T$ years at an annual rate\ix{Annual rate} $r$. Under \emph{simple interest}\index{Simple interest} the
interest is paid on the principal only, so the terminal value is $P(1+rT)$. Under \emph{compound
interest} with $m$ compounding\ix{Compounding} periods per year, interest earns interest at the end of each period:
\begin{equation}\label{eq:periodic}
  V_T = P\left(1+\frac{r}{m}\right)^{mT}.
\end{equation}
The rate $r$ in \cref{eq:periodic} is a \emph{nominal} or quoted rate\ix{Quoted rate}. The \emph{effective annual rate\ix{Effective annual rate}}\index{Effective annual rate}
is the growth over one year, $r_{\mathrm{eff}} = (1+r/m)^m - 1$. For a positive quoted rate, the case
used throughout, the effective rate rises as $m$ grows, but it does not grow without bound.

\begin{proposition}\label{prop:limit}
For any real $r$, $\lim_{m\to\infty}(1+r/m)^m = e^{r}$.
\end{proposition}
\begin{proof}
For $m>\lvert r\rvert$ the base $1+r/m$ is positive, so $(1+r/m)^m = \exp\{m\ln(1+r/m)\}$. The Taylor expansion $\ln(1+x) = x - x^2/2 + O(x^3)$ with
$x=r/m$ gives $m\ln(1+r/m) = r - r^2/(2m) + O(m^{-2}) \to r$, and the exponential is continuous.
\end{proof}

The limit defines \emph{continuous compounding\ix{Continuous compounding}}\index{Continuous compounding}: $V_T = Pe^{rT}$. Continuous rates are the natural
language of derivatives pricing and of the stochastic models in later chapters. Converting between
conventions only requires that the two give the same growth over a year: a continuously compounded
rate $r_c$ and an annual effective rate $r_a$ are equivalent when $e^{r_c} = 1+r_a$, so
$r_c = \ln(1+r_a)$.

\begin{example}{Effective rates of a quoted 8\% rate}\label{ex:ear}
A quoted rate of 8\% gives an effective annual rate of \earA\% with annual compounding, \earB\% with
semiannual, \earC\% with quarterly, \earD\% with monthly and \earE\% with daily (365 periods)
compounding. The continuous limit is $e^{0.08}-1 = \earF\%$. Going the other way, an annual
effective rate of 8\% corresponds to a continuously compounded rate of $\ln 1.08 = \ccEquivEight\%$.
Panel (a) of \cref{fig:compounding} shows how quickly the sequence settles: most of the gap between
annual and continuous compounding is closed by quarterly compounding.
\end{example}

\begin{figure}[t]
  \centering
  \includegraphics[width=\textwidth]{ch01/fig01_compounding}
  \caption{(a) Effective annual rate of a nominal 8\% rate as the number of compounding periods per
  year grows, with the continuous limit $e^{0.08}-1$. (b) The log return $\ln(1+R)$ against the simple
  return $R$. The two agree near zero and diverge for large moves; the log return can never exceed the
  simple return because $\ln(1+R)\le R$.}
  \label{fig:compounding}
\end{figure}

\subsection{Returns that add over time and returns that add across assets}

Let a position be worth $V_{t-1}>0$ at the start of a period and $V_t$ at the end, and let $D_t$ be
any cash it pays the holder during the period. The \emph{simple} or total return is
$R_t = (V_t+D_t-V_{t-1})/V_{t-1}$ and the \emph{log return\ix{Log return}}\index{Log return}, or continuously compounded\ix{Compounding} return, is
$r_t = \ln(1+R_t)$. The logarithm needs $1+R_t>0$, so a total loss has no finite log return. Here $D_t$ is treated as paid at the end of the period; a payment made part way
through is handled by splitting the period at the payment date. If the cash is reinvested in the
position, a holder's wealth $W_t$ grows by the factor $1+R_t$ each period. Each kind of return has one property the other lacks.

\begin{proposition}\label{prop:additivity}
(i) Log returns add over time: if $W_t = W_{t-1}(1+R_t)$ with every $W_t>0$, the log return from $0$
to $T$, $\ln(W_T/W_0)$, equals $\sum_{t=1}^T r_t$.
(ii) Simple returns\ix{Simple return} add across assets: if a portfolio holds value weights $w_i$ (summing to one) at
the start of the period and no money is added or withdrawn during it, its simple return is
$R_p = \sum_i w_i R_i$.
\end{proposition}
\begin{proof}
(i) $\ln(W_T/W_0) = \ln\prod_{t}(W_t/W_{t-1}) = \sum_t \ln(1+R_t) = \sum_t r_t$. (ii) Let the portfolio
be worth $V_{p,t-1}$, so that it holds $w_iV_{p,t-1}$ in asset $i$. With no money added or withdrawn,
each holding grows to $w_iV_{p,t-1}(1+R_i)$, distributions included, so
$V_{p,t} = \sum_i w_i V_{p,t-1}(1+R_i)$ and $V_{p,t}/V_{p,t-1}-1 = \sum_i w_i R_i$.
\end{proof}

Neither property carries over to the other definition. The log return of a portfolio is not the
weighted average\ix{Weighted average} of the constituents' log returns, and simple returns compound rather than add. Log
returns make compounding through time simple, which is why they are used for multi-period growth,
volatility scaling and the stochastic models of Chapters~3 and~5. Simple returns aggregate exactly
across assets over a common period, which is why they are used for portfolio returns and
attribution. Neither choice removes the need to account for cash flows or for serial dependence.

\begin{example}{Up 50\%, down a third}\label{ex:updown}
A fund unit goes from 100 to 150 and back to 100. The simple returns are $+50\%$ and $-33.33\%$, with
an arithmetic mean\ix{Arithmetic mean} of $+8.33\%$ per period, although the investor ends where they began. The log
returns are $\ln 1.5 = \upLog\%$ and $\ln(100/150) = \downLog\%$, which sum to zero. The arithmetic
mean of simple returns overstates the growth rate\ix{Growth rate} whenever returns vary. When returns are small,
the mean log return is close to the arithmetic mean of the simple returns less half their variance,
and the compound growth rate\ix{Compound growth rate}, the exponential of the mean log return less one, is close to the mean
log return itself. Here,
with the variance computed by dividing by the number of periods, that gives
$\udMean\%-\udHalfVar\%=\udApprox\%$ against a true growth rate of 0\%. With the sample variance, which divides by one less, it gives
\udApproxSample\%. With only two large, unequal returns the second-order approximation is unreliable, and changing
the variance convention changes its result substantially; the expansion is best read as a statement
about mean log growth over many periods of small returns.
The gap matters for volatile strategies and for any fee or hurdle computed on arithmetic averages.
\end{example}

\FloatBarrier
% =====================================================================================
\section{Measuring a fund that calls and returns capital}
\label{sec:irr}

\subsection{The internal rate of return}

Private funds call capital from investors when deals are found and return it when deals are sold.
The investor does not control the timing, so a return measure for the fund should reflect the actual
dates and sizes of the flows. The \emph{internal rate of return\ix{Internal rate of return}\ix{Rate of return}}\index{Internal rate of return} does this, though on its own it
does not say whether the fund did better than a public alternative; \cref{sec:multiples} returns to
that question.

\begin{definition}\label{def:irr}
For cash flows $\mathrm{CF}_0,\dots,\mathrm{CF}_T$ at equally spaced annual dates $0,\dots,T$ (positive when
received by the investor), the internal rate of return is any rate $\IRR > -1$ that solves
\begin{equation}\label{eq:irr}
  \NPV(\IRR) = \sum_{t=0}^{T}\frac{\mathrm{CF}_t}{(1+\IRR)^{t}} = 0 .
\end{equation}
\end{definition}
For flows on irregular dates $t_j$, measured in years from the first flow, the discount factor
becomes $(1+\IRR)^{-t_j}$.

The IRR is a money-weighted return: it depends jointly on the amounts and dates of all contributions,
distributions and the terminal value. For
a fund that has not finished, the \emph{since-inception} or \emph{interim} IRR treats the current net
asset value\ix{Net asset value} (NAV) as if it were a final distribution received today. That makes the interim IRR only
as good as the NAV, a point that drives the J-curve\index{J-curve} in \cref{sec:jcurve}.

\subsection{Fund A}

Fund A, used from here to \cref{sec:multiples}, is small enough to check with a calculator and
still shows the J-curve, the fee drag\ix{Fee drag} and the IRR's blind spots. Fund A buys one deal a year for
five years, at costs of 20, 22, 20, 15 and 10. Each deal gains 15\% a year in economic value\ix{Economic value} and is
sold after five years. The fund charges a fee of 2 a year for nine years; in years 0 to 4 the fee is
called from investors\ix{Investor} together with the deal cost, and in years 5 to 8 it is deducted from
distributions. Reporting is conservative: a deal is carried at 90\% of cost in the year it is bought,
at cost in its next two years, and at economic value from its third anniversary until it is sold.
\Cref{tab:funda} lists the flows to investors and the reported NAV after each date's flows.

\begin{table}[t]
  \centering
  \caption{Fund A: contributions, distributions and reported NAV (end of year, after flows).}
  \label{tab:funda}
  \small\setlength{\tabcolsep}{4.5pt}
  \begin{tabular}{@{}lrrrrrrrrrr@{}}
    \toprule
    Year & 0 & 1 & 2 & 3 & 4 & 5 & 6 & 7 & 8 & 9\\
    \midrule
    Contributions & \fcA & \fcB & \fcC & \fcD & \fcE & \fcF & \fcG & \fcH & \fcI & \fcJ\\
    Distributions & \fdA & \fdB & \fdC & \fdD & \fdE & \fdF & \fdG & \fdH & \fdI & \fdJ\\
    Reported NAV  & \fnA & \fnB & \fnC & \fnD & \fnE & \fnF & \fnG & \fnH & \fnI & \fnJ\\
    \bottomrule
  \end{tabular}
\end{table}

Investors pay in \fundPaidIn{} and receive \fundDistrib. Solving \cref{eq:irr} for the net flows
(distributions minus contributions) gives a final IRR of \fundIRR\%, below the 15\% earned by each deal
because of the management fee. With the fee set to zero, every deal's flows and the fund's flows
have an IRR of exactly \fundIRRNoFee\%, so the whole gap of \feeDrag{} percentage points\ix{Percentage point} is the drag
of the fee. No formula in radicals solves every polynomial of degree five or more, and the fund's NPV\ix{Net present value} is a polynomial of degree nine in $1/(1+r)$, so the IRR is found
numerically\ix{Numerical methods}. Newton's method\ix{Newton's method}, $r_{k+1} = r_k - \NPV(r_k)/\NPV'(r_k)$,
usually converges in a handful of steps from a sensible start, but no start guarantees it. The
listing raises an error if the derivative is zero or a Newton step reaches or crosses $-100\%$. It
stops when a step is small, which does not by itself show that the NPV is small at the answer:
\Needspace*{21\baselineskip}
\lstinputlisting[language=Python]{code/ch01/snippets/irr_newton.py}
Run as shown, on the two-decimal flows of \cref{tab:funda}, the listing prints 11.5298\%; the
unrounded flows give 11.5294\%, and both round to \fundIRR\%. Production code is better off bracketing the root and using Brent's method\ix{Root finding}, as
\texttt{ch01lib.irr\_all} does, because Newton's method can also converge to the wrong root when there
is more than one. The \texttt{irr} function of \texttt{numpy-financial} \parencite{numpyfinancial}
returns the same value.

\FloatBarrier
% =====================================================================================
\section{Why young funds look bad}
\label{sec:jcurve}

In the early years of a drawdown fund the interim IRR\ix{Interim IRR} is often negative even when the fund will end
up doing well. Fees are charged on committed capital from the first day, before investments have had time to grow.
Some transaction and organisational costs are expensed at once; others are capitalised into the cost
of an investment, depending on the fund's accounting policy. Valuation\ix{Valuation} adds to the effect. Fair value
estimates can rise or fall before an asset is sold, depending on the evidence available and the
fund's valuation policy, and in some funds the evidence needed to justify a higher value arrives later than the evidence for
an impairment. Fund A's rule of 90\% of cost in the year of
purchase and cost for two more years is a stylised version of that caution.

\Cref{fig:jcurve} shows the result for Fund A. Cumulative net cash\ix{Cumulative net cash flow} flow to investors falls to
\cumTrough{} at year \cumTroughYear{} and turns positive in year \breakEvenYear. The interim IRR starts
at \interimOne\% after one year and \interimTwo\% after two, is close to zero (\interimThree\%) after
three, and reaches \interimFive\% after five. From year 7 it stays within ten basis points of the
final \fundIRR\%, because by then every deal still held is carried at its economic value. The plotted
interim IRRs trace the letter J. Early negative interim returns say little about the manager; the fees, the deals and the basis of the
valuations say more. In Fund A, whose deals grow steadily in the model, they reflect fees and the assumed valuation
schedule. A fund that shows a strongly positive
interim IRR in its first two years should prompt questions about how its assets are being marked.

\begin{figure}[t]
  \centering
  \includegraphics[width=\textwidth]{ch01/fig04_jcurve}
  \caption{The J-curve of Fund A. (a) Contributions, distributions and cumulative net cash flow to
  investors. (b) Since-inception IRR using the reported NAV at each year end.}
  \label{fig:jcurve}
\end{figure}

\FloatBarrier
% =====================================================================================
\section{When the IRR misleads}
\label{sec:irrproblems}

\subsection{Several roots, or none}

Write $x = 1/(1+r)$. Then \cref{eq:irr} is the polynomial equation $\sum_t \mathrm{CF}_t x^t = 0$, and
Descartes\ix{Sign rule}' rule of signs bounds the number of positive roots $x$ (hence the number of IRRs\ix{Internal rate of return} above
$-100\%$) by the number of sign changes in the sequence of cash flows\ix{Cash flow}. Descartes' rule gives only an upper
bound. For the commonest pattern there is a sharper statement with a direct proof.

\begin{definition}\label{def:conventional}
Cash flows $\mathrm{CF}_0,\dots,\mathrm{CF}_T$ are \emph{conventional} if there is a date $k<T$ such that
$\mathrm{CF}_t\le 0$ for $t\le k$, $\mathrm{CF}_t\ge 0$ for $t>k$, $\mathrm{CF}_t<0$ for at least one
$t\le k$ and $\mathrm{CF}_t>0$ for at least one $t>k$.
\end{definition}

\begin{proposition}\label{prop:uniqueirr}
Conventional cash flows have exactly one IRR in $(-1,\infty)$. The IRR is positive if
$\sum_t\mathrm{CF}_t>0$, zero if $\sum_t\mathrm{CF}_t=0$ and negative if $\sum_t\mathrm{CF}_t<0$.
In addition $\NPV(r)>0$ for every $r$ below the IRR and $\NPV(r)<0$ for every $r$ above it.
\end{proposition}

\begin{proof}
Since $(1+r)^k>0$ for $r>-1$, $\NPV(r)$ and
\begin{equation}\label{eq:fk}
  f(r) = (1+r)^k\,\NPV(r) = \sum_{t\le k}\mathrm{CF}_t\,(1+r)^{k-t} + \sum_{t>k}\mathrm{CF}_t\,(1+r)^{-(t-k)}
\end{equation}
have the same zeros and the same sign. In the first sum each exponent $k-t$ is non-negative and each
$\mathrm{CF}_t$ is non-positive; in the second each exponent is negative and each $\mathrm{CF}_t$ is
non-negative. Differentiating term by term,
\begin{equation*}
  f'(r) = \sum_{t<k}(k-t)\,\mathrm{CF}_t\,(1+r)^{k-t-1} - \sum_{t>k}(t-k)\,\mathrm{CF}_t\,(1+r)^{-(t-k)-1} .
\end{equation*}
Every term of the first sum is $\le 0$ and every term of the second is $\ge 0$ and subtracted, so
$f'(r)\le 0$. At least one inflow is strictly positive and enters the second sum with $t-k\ge 1$, so
$f'(r)<0$ and $f$ is strictly decreasing on $(-1,\infty)$. Next, the limits. As $r\downarrow -1$ the
terms with $t>k$ grow without bound, because $(1+r)^{-(t-k)}\to\infty$ and at least one of them has a
positive coefficient, while the first sum stays bounded; hence $f(r)\to+\infty$. As $r\to\infty$ the
second sum tends to zero and the first tends to $-\infty$ if some outflow occurs before $k$, or to
$\mathrm{CF}_k$ if the only outflow is at $k$. In both cases $\lim_{r\to\infty}f(r)<0$, because at least one
outflow is strictly negative. A continuous, strictly decreasing function that is positive near $-1$ and
negative for large $r$ has exactly one zero $r^\ast$, with $f>0$ to its left and $f<0$ to its right. Finally
$f(0)=\NPV(0)=\sum_t\mathrm{CF}_t$, and since $f$ is strictly decreasing, $f(0)>0$ holds if and only if
$0<r^\ast$, which gives the sign of the IRR.
\end{proof}

The last statement of \cref{prop:uniqueirr} is what makes the IRR rule work for conventional
projects: accepting when the IRR exceeds the cost of capital $r$ is the same as accepting when
$\NPV(r)>0$. Fund A is of this kind, and panel (a) of \cref{fig:npvprofiles} shows its single crossing.

When cash flows change sign more than once, the NPV\ix{Net present value} equation can have several IRRs or none. Consider $-100$ now, $+230$ in a year
and $-132$ in two years, a pattern that arises when a fund recalls distributions or a project has
clean-up costs. Setting $-100 + 230x - 132x^2 = 0$ gives $x = (230\pm 10)/264$, so the IRR is both
\multiLo\% and \multiHi\% (panel (b)). Both solve the defining equation, so the IRR cannot say on
its own whether the project is worth doing; the NPV is positive only for discount rates\ix{Discount rate} between the
two. Now consider $+100$, $-300$, $+250$, a
borrowing-type pattern in which money is received first. The quadratic $100 - 300x + 250x^2$ has a
negative discriminant, $300^2 - 4\cdot 100\cdot 250 < 0$, so there is no real IRR at all; the NPV is
positive at every discount rate, with a minimum value of \noneMinNPV{} (panel (c)).

\begin{figure}[t]
  \centering
  \includegraphics[width=\textwidth]{ch01/fig02_npv_profiles}
  \caption{NPV profiles. (a) Fund A has one IRR at \fundIRR\%. (b) Flows that change sign twice can
  have two IRRs, here \multiLo\% and \multiHi\%. (c) A borrowing-type pattern can have none: the NPV
  stays positive at every rate.}
  \label{fig:npvprofiles}
\end{figure}

\subsection{Scale and timing}

The IRR is a rate and says nothing about the amount of wealth created. A deal that turns 1 into 2 in a
year has an IRR of 100\%, while one that turns 1{,}000 into 1{,}500 has 50\%; if only one can be
done, the first has the higher rate but creates far less wealth at any ordinary cost of capital\ix{Cost of capital}. Timing creates a subtler problem. Compare a short project, $-100$ now and $+130$ in a year,
with a long project, $-100$ now and $+200$ in three years. Their IRRs are \irrShort\% and \irrLong\%, so
IRR ranks the short project first. Their NPVs\ix{Net present value} are equal where
\begin{equation*}
  -100 + \frac{130}{1+r} = -100 + \frac{200}{(1+r)^3}
  \iff (1+r)^2 = \frac{200}{130}
  \iff r = \sqrt{200/130}-1 = \crossRate\% ,
\end{equation*}
the \emph{crossover rate}\index{Crossover rate}. At any cost of capital below \crossRate\% the long project creates more
value; at 10\% its NPV is \npvLongTenAbs{} against \npvShortTenAbs{} for the short one
(\cref{fig:crossover}). The decision rule here is NPV at the cost of capital, and it needs no
assumption about what happens to the 130. Reinvestment enters only if we ask a different question,
about wealth at year 3: the short project ends with more than 200 then if its 130 can be put to work
for two more years at above \crossRate\%, since $130(1+r)^2>200$ exactly when $r>\crossRate\%$. The IRR alone cannot rank mutually exclusive projects
whose cash flows fall on different dates.

\begin{figure}[t]
  \centering
  \includegraphics[width=0.86\textwidth]{ch01/fig03_crossover}
  \caption{The timing problem. The short project has the higher IRR, but the long project has the
  higher NPV at every discount rate below the crossover rate of \crossRate\%.}
  \label{fig:crossover}
\end{figure}

\subsection{IRRs do not average}

Fund X pays $-100$ then $+150$ a year later (IRR\ix{Internal rate of return} \irrX\%); Fund Y pays $-100$ then
$+121$ two years later (IRR \irrY\%). A portfolio holding both has flows $-200$, $+150$, $+121$ and an
IRR of \irrXY\%, not the simple average of 30\%. The combined IRR solves one NPV equation for all the dated flows, so it is no fixed-weight average
of the two IRRs; Fund Y keeps its capital at work for longer and pulls the result towards its own
rate. This matters whenever IRRs
of funds, vintages\ix{Vintage year} or deals are averaged in a pitch book or a database.

\subsection{Reinvestment and the IRR}

Computing an IRR\ix{Internal rate of return} assumes nothing about what an investor does with a distribution: it is a rate at
which the net present value of the flows is zero. Reinvestment\ix{Reinvestment} enters only when the IRR is read as
the rate at which the investor's total wealth grew. That reading needs an assumption about money
outside the fund: it holds when every interim distribution is reinvested at the IRR and every later
contribution would otherwise have earned the IRR, because at that rate the compounded\ix{Compounding} inflows exactly
match the compounded outflows. When a fund with a high IRR makes a large early distribution that the investor can only reinvest
at a lower rate, the reading overstates the growth of the investor's wealth. The same arithmetic lets
a quick early exit lift a fund's IRR far more than its multiple \parencite{phalippou2008}. The modified
IRR makes the reinvestment rate explicit.

\FloatBarrier
% =====================================================================================
\section{Explicit rates, and whose timing counts}
\label{sec:mirr}

\subsection{The modified IRR}

The \emph{modified internal rate of return\ix{Internal rate of return}}\index{Modified internal rate of return} builds one annualised rate from explicit financing and
reinvestment\ix{Reinvestment} rates: a finance rate $f>-1$ at which outflows are discounted to date $0$, and a
reinvestment rate $\rho>-1$ at which inflows are compounded to date $T>0$. With $\mathrm{CF}_t^- = \min(\mathrm{CF}_t,0)$
and $\mathrm{CF}_t^+ = \max(\mathrm{CF}_t,0)$,
\begin{equation}\label{eq:mirr}
  \MIRR = \left(\frac{\sum_t \mathrm{CF}_t^+(1+\rho)^{T-t}}{-\sum_t \mathrm{CF}_t^-(1+f)^{-t}}\right)^{1/T}-1 .
\end{equation}
Because the right-hand side is an explicit formula, the MIRR\ix{Modified internal rate of return} always exists and is unique when there is
at least one outflow and one inflow. Its weakness is that the answer depends on the two rates the
analyst chooses. When both rates equal the cost of capital\ix{Cost of capital} the MIRR gives the same accept or reject
decision as the NPV\ix{Net present value}, whatever the pattern of signs.

\begin{proposition}\label{prop:mirr}
Let the flows contain at least one outflow and one inflow, and set $f=\rho=r>-1$. Write
$\mathrm{PV}^- = -\sum_t\mathrm{CF}_t^-(1+r)^{-t}$ and $\mathrm{FV}^+ = \sum_t\mathrm{CF}_t^+(1+r)^{T-t}$.
Then
\begin{equation}\label{eq:mirrnpv}
  \NPV(r) = \mathrm{PV}^-\Bigl[\Bigl(\frac{1+\MIRR}{1+r}\Bigr)^{T}-1\Bigr],
\end{equation}
so $\MIRR>r$, $\MIRR=r$ or $\MIRR<r$ exactly when $\NPV(r)>0$, $=0$ or $<0$.
\end{proposition}

\begin{proof}
Split each flow into its positive and negative parts, $\mathrm{CF}_t = \mathrm{CF}_t^+ + \mathrm{CF}_t^-$.
Then
\begin{align*}
  \NPV(r) &= \sum_t \mathrm{CF}_t^+(1+r)^{-t} + \sum_t \mathrm{CF}_t^-(1+r)^{-t}
           = \frac{\mathrm{FV}^+}{(1+r)^T} - \mathrm{PV}^- .
\end{align*}
By \cref{eq:mirr} with $f=\rho=r$, $(1+\MIRR)^T = \mathrm{FV}^+/\mathrm{PV}^-$. Substituting
$\mathrm{FV}^+ = \mathrm{PV}^-(1+\MIRR)^T$ gives \cref{eq:mirrnpv}. Since $\mathrm{PV}^->0$ and
$x\mapsto x^T$ is strictly increasing on $(0,\infty)$, the sign of $\NPV(r)$ is the sign of
$(1+\MIRR)/(1+r)-1$.
\end{proof}

\begin{example}{MIRR of the two-root project}\label{ex:mirr}
Take $-100$, $+230$, $-132$ with a finance rate of 8\% and a reinvestment rate of 12\%. The present value
of the outflows is $100 + 132/1.08^2 = \mirrPVneg$. The future value of the inflow at year 2 is
$230\times 1.12 = \mirrFVpos$. The MIRR is $(\mirrFVpos/\mirrPVneg)^{1/2}-1 = \mirrMulti\%$, a single
answer where the IRR gave two. For Fund A, the MIRR is \mirrFundEight\% when both rates are 8\% and
\mirrFundTwelve\% when distributions are reinvested at 12\% and outflows still financed at 8\%. Both
are below the IRR of \fundIRR\%. In the second case the reinvestment rate is above the IRR, but
financing the outflows at 8\%, well below it, pulls the MIRR down by more than reinvesting at 12\%
lifts it. With both rates at 12\% the MIRR is \mirrFundTwelveTwelve\%, above the IRR. For conventional flows an MIRR computed at a single rate $r$ lies between $r$ and the IRR. If $r$ is
below the IRR, $\NPV(r)>0$ by \cref{prop:uniqueirr}, so \cref{prop:mirr} puts the MIRR above $r$. Since
$(1+\MIRR)^T=\mathrm{FV}^+/\mathrm{PV}^-$ never falls as $r$ rises and equals $(1+\IRR)^T$ at $r=\IRR$,
the MIRR cannot exceed the IRR; for $r$ below the IRR it equals the IRR only when every outflow falls
at date $0$ and every inflow at date $T$. Above the IRR both inequalities reverse. Since the MIRR at 8\% exceeds 8\%, \cref{prop:mirr} confirms that Fund A has a positive NPV at an 8\%
cost of capital, as \cref{prop:uniqueirr} also implies from its IRR.
\end{example}

\subsection{Whose timing counts}

A \emph{time-weighted\ix{Time-weighted return} return}\index{Time-weighted return} links the returns of successive sub-periods, $\prod_k(1+R_k)-1$, with
sub-periods cut at every external flow. It measures what one unit invested at the start would have
earned and is therefore independent of when the investor added or withdrew money. A
\emph{money-weighted\ix{Money-weighted return} return}\index{Money-weighted return} is the IRR of the investor's own flows and does depend on that timing.
The independence assumes that the position can be valued reliably on each flow date, which is easy
for a listed fund and hard for an appraised one.
Time weighting\ix{Weighting scheme} is the right way to judge a manager who does not control flows, such as a mutual fund
manager\ix{Fund manager}. Money weighting is especially informative when the manager controls the timing and amount of
capital deployed, as the general partner of a drawdown fund does, although the result still depends
on when the cash moved and on how the remaining assets are valued.

\begin{example}{The same fund, two answers}\label{ex:twr}
An investor puts 100 into a fund. The fund gains 20\% in the first period, the investor adds 200 when
the account is worth 120, and the fund then loses 10\%, leaving 288. The time-weighted return is
$1.2\times 0.9 - 1 = \twrTotal\%$ over two periods, or \twrPer\% per period. The money-weighted return
solves $-100(1+r)^2 - 200(1+r) + 288 = 0$ and equals \dwrPer\% per period. The fund's own performance was
positive; the investor's experience was negative because most of the money arrived just before the
loss.
\end{example}

\FloatBarrier
% =====================================================================================
\section{Multiples and the public market yardstick}
\label{sec:multiples}

Private fund investors usually read three multiples\ix{Comparable companies}\ix{Valuation multiple} next to the IRR. With cumulative paid-in capital\ix{Paid-in capital}
$\mathrm{PIC}$, cumulative distributions $D$ and current NAV\ix{Net asset value},
\begin{equation}\label{eq:multiples}
  \mathrm{DPI} = \frac{D}{\mathrm{PIC}},\qquad
  \mathrm{RVPI} = \frac{\mathrm{NAV}}{\mathrm{PIC}},\qquad
  \mathrm{TVPI} = \mathrm{DPI} + \mathrm{RVPI} .
\end{equation}
DPI\ix{Distributed to paid-in} is the realised part of performance, RVPI\ix{Residual value to paid-in} the unrealised part, and TVPI\ix{Total value to paid-in} the total. Multiples
ignore timing, so a quick early exit cannot inflate them, and they do not penalise slow
distributions even though waiting lowers their present value\ix{Present value}. For Fund A at the end of year 5, DPI is \dpiFive, RVPI is \rvpiFive{} and TVPI is \tvpiFive. At
the end of its life, all value has been distributed and DPI equals TVPI at \tvpiEnd.

Neither the IRR nor the multiples say whether the fund beat what the same money could have earned in
public markets\ix{Public markets}. A \emph{public market equivalent}\index{Public market equivalent} (PME\ix{Kaplan-Schoar PME}) answers that by pricing the fund's flows with
a public index. Let $I_t$ be the index level\ix{Index level} at date $t$, $C_t$ the contributions and $D_t$ the
distributions. The PME of \textcite{kaplan2005} carries every flow forward to date $T$
with the index return:
\begin{equation}\label{eq:kspme}
  \mathrm{PME}_{\mathrm{KS}} = \frac{\sum_t D_t\,(I_T/I_t) + \mathrm{NAV}_T}{\sum_t C_t\,(I_T/I_t)} .
\end{equation}
A value above one means the fund delivered more than an index investment making the same dated
contributions and distributions would have, provided the index is a fair, investable comparison for
the fund's assets.
\Cref{prop:pme} below states this precisely.

\begin{proposition}\label{prop:pme}
Let $\mathrm{NPV}_I = \sum_t (D_t - C_t)/I_t + \mathrm{NAV}_T/I_T$ be the net present value\ix{Net present value} of the fund's
flows when each date is discounted by the index level, and suppose $\sum_t C_t>0$. Then
\begin{equation}\label{eq:pmenpv}
  \mathrm{PME}_{\mathrm{KS}} - 1 = \frac{\mathrm{NPV}_I}{\sum_t C_t/I_t},
\end{equation}
so $\mathrm{PME}_{\mathrm{KS}}>1$ if and only if $\mathrm{NPV}_I>0$. If the index grows at a constant rate
$g$, so that $I_t = I_0(1+g)^t$, then $\mathrm{PME}_{\mathrm{KS}}>1$ if and only if $\NPV(g)>0$ for the net
flows with NAV treated as a final distribution, and for conventional flows this holds if and only if
$\IRR>g$.
\end{proposition}

\begin{proof}
Divide the numerator and denominator of \cref{eq:kspme} by $I_T>0$:
\begin{equation*}
  \mathrm{PME}_{\mathrm{KS}} = \frac{\sum_t D_t/I_t + \mathrm{NAV}_T/I_T}{\sum_t C_t/I_t},
  \qquad\text{so}\qquad
  \mathrm{PME}_{\mathrm{KS}} - 1 = \frac{\sum_t (D_t-C_t)/I_t + \mathrm{NAV}_T/I_T}{\sum_t C_t/I_t},
\end{equation*}
which is \cref{eq:pmenpv}. The denominator is positive, so the sign of $\mathrm{PME}_{\mathrm{KS}}-1$ is the
sign of $\mathrm{NPV}_I$. With $I_t = I_0(1+g)^t$, $\mathrm{NPV}_I = \NPV(g)/I_0$, which has the sign of
$\NPV(g)$. For conventional flows \cref{prop:uniqueirr} says that $\NPV(g)>0$ exactly when $g$ lies below
the IRR.
\end{proof}

The earlier method of \textcite{long1996} builds a shadow portfolio\ix{Shadow portfolio}: each contribution buys
the index, each distribution sells the same amount of it, and the IRR of the fund's flows with the
shadow portfolio's final value is compared with the fund's IRR. A distribution larger than the shadow
portfolio's value on the day it is paid drives the shadow holding below zero. The shadow portfolio is then short the index, or borrowing against it,
which is a different investment from the long-only alternative the method sets out to describe.

\begin{table}[t]
  \centering
  \caption{Hypothetical public index used for the PME calculations (level at the end of each year).}
  \label{tab:index}
  \small\setlength{\tabcolsep}{4.5pt}
  \begin{tabular}{@{}lrrrrrrrrrr@{}}
    \toprule
    Year & 0 & 1 & 2 & 3 & 4 & 5 & 6 & 7 & 8 & 9\\
    \midrule
    Return (\%) &  & \irB & \irC & \irD & \irE & \irF & \irG & \irH & \irI & \irJ\\
    Index level & \ixA & \ixB & \ixC & \ixD & \ixE & \ixF & \ixG & \ixH & \ixI & \ixJ\\
    \bottomrule
  \end{tabular}
\end{table}

With the hypothetical index of \cref{tab:index}, which compounds at \indexCAGR\% a year, Fund A has a
Kaplan and Schoar PME of \ksPME: the index-adjusted value of its distributions and final NAV is
\ksPME{} times the index-adjusted value of its contributions.
The Long and Nickels\ix{Long-Nickels PME} shadow portfolio shows a known flaw (\cref{fig:pme}). Because the fund distributes
far more than the index portfolio has grown to, the shadow portfolio is sold down below zero and ends
at \lnValue, the negative holding described above. The last shadow flow
is a net outflow, so the flows change sign twice, and their NPV is zero at about 5.8\% and again near
$-56\%$: the Long and Nickels IRR is not even unique. Ratio methods such as \cref{eq:kspme} do not have this problem.

\begin{figure}[t]
  \centering
  \includegraphics[width=0.86\textwidth]{ch01/fig05_pme}
  \caption{Fund A's NAV against the Long and Nickels shadow index portfolio built from the same
  contributions and distributions. Strong distributions drive the shadow portfolio below zero.}
  \label{fig:pme}
\end{figure}

\FloatBarrier
% =====================================================================================
\section{Returns on derivatives with no entry price}
\label{sec:notional}

Forwards, futures\ix{Futures contract} and swaps usually have zero market value when they are entered. Margin may have to
be posted, but that money is collateral and is not the position's price, and any fees or upfront payments are a separate
matter. A return defined as profit divided by initial value is then undefined, yet managed futures
funds and overlay programmes must report returns. The standard solution is to express the position relative to its \emph{notional
principal\ix{Face value}}, the face amount of the underlying exposure\ix{Exposure}, and to specify how much cash backs it.

A position is \emph{fully collateralised} when cash equal to the notional principal is set aside and
earns a short-term rate. For a suitable futures contract the position then approximates the economic exposure of holding the
underlying, though its performance can differ: futures pricing, the contract's specification and
the return on the collateral all open gaps. Over a period its profit is the change in value of the derivative plus the interest on
the collateral. Writing $F_t$ for the futures price, $N$ for the notional and $y$ for the
collateral yield over the period, and assuming a fixed exposure, collateral that earns that yield and no
financing or transaction costs\ix{Transaction costs}, the simple return\ix{Simple return} on a fully collateralised long futures position is
\begin{equation}\label{eq:fullcoll}
  R^{\text{full}} = \frac{N\,(F_t-F_{t-1})/F_{t-1} + yN}{N} = \frac{F_t-F_{t-1}}{F_{t-1}} + y .
\end{equation}
If only a fraction $c$ of the notional is held as collateral, the same futures gain is earned on a
smaller capital base, while the interest shrinks in proportion. The \emph{partially collateralised} return is
\begin{equation}\label{eq:partcoll}
  R^{\text{part}} = \frac{N\,(F_t-F_{t-1})/F_{t-1} + y\,cN}{cN}
  = \frac{1}{c}\,\frac{F_t-F_{t-1}}{F_{t-1}} + y ,
\end{equation}
so the leverage\ix{Leverage} factor is $1/c$. The log return\ix{Log return} in either case is $\ln(1+R)$, the form to use when returns are compounded or summed
through time.

\begin{example}{A futures position at two collateral levels}\label{ex:notional}
A manager buys equity index futures with a notional principal of \$10{,}000{,}000. Over a month the
futures price rises 2\% and Treasury bills\ix{Treasury bill} yield 0.40\% for the month. Fully collateralised, the profit
is \$200{,}000 from the futures plus \$40{,}000 interest, \$\pnlFull{} in total, a simple return of
\simpleFull\% and a log return of $\ln(1.024) = \logFull\%$. With 25\% collateral (\$2{,}500{,}000) the
interest falls to \$10{,}000, the profit is \$\pnlPart{}, and the simple return on capital is
\simplePart\% (log return \logPart\%). If the futures had instead fallen 2\%, the partially
collateralised return would have been \simplePartLoss\% (log return \logPartLoss\%). Leverage of four
scales gains and losses alike. In both cases the log return lies below the simple return, so in log
terms a gain looks smaller and a loss looks larger.
\end{example}

\FloatBarrier
% =====================================================================================
\section{Models and track records}
\label{sec:distributions}

A model describes the returns that might occur; a track record\ix{Track record} describes the returns that did.
Twenty years of a buyout fund's quarterly NAVs is eighty appraisals\ix{Appraisal}, not eighty prices. The
empirical distribution of a track record, whether shown as a histogram, a density estimate or a
table of summary statistics, can be used to estimate a model's parameters once the model has been
chosen, and the quality of those estimates depends
on how much data there is and on whether the past is a fair sample of the future. For private
assets the track record is itself built from estimates.

\subsection{The normal and lognormal laws}

The normal distribution\ix{Normal distribution} $\N(\mu,\sigma^2)$ has density\ix{Density}
$(2\pi\sigma^2)^{-1/2}\exp\{-(x-\mu)^2/(2\sigma^2)\}$. It is a useful benchmark for log returns\ix{Log return} for two
reasons. First, the central limit theorem says that the standardised sum of independent, identically distributed\ix{Identically distributed}
shocks with finite variance\ix{Finite variance} tends to a normal law, and versions of it hold for dependent shocks under
further conditions. Second, linear combinations of jointly normal\ix{Gaussian process} variables are normal, so a portfolio of assets
with jointly normal returns has a normal return. The first reason is the
weaker one for alternatives: jumps, long swings in volatility and heavy tails can make the approach to
the normal slow, so that a sum is still far from normal at the horizons an investor sees. One
flaw as a model of simple returns\ix{Simple return} is obvious: it gives positive probability to losing more than
everything.

The lognormal\ix{Lognormal distribution}\index{Lognormal distribution} model fixes that flaw. If the log return $Y = \ln(1+R)$ is $\N(\mu,\sigma^2)$, then the
gross return $X = 1+R = e^Y$ is lognormal and always positive. Its moments follow from the moment
generating function of the normal.

\begin{proposition}\label{prop:lnmoments}
Let $X=e^Y$ with $Y\sim\N(\mu,\sigma^2)$, $\sigma>0$, and write $w=e^{\sigma^2}$ and $m=e^{\mu+\sigma^2/2}$. Then
\begin{align}
  \E[X] &= m, \qquad \Var[X] = m^2(w-1), \label{eq:lnmoments}\\
  S &= (w+2)\sqrt{w-1}, \qquad K_e = w^4+2w^3+3w^2-6, \label{eq:lnshape}
\end{align}
the median\ix{Median} of $X$ is $e^\mu$ and its mode is $e^{\mu-\sigma^2}$, so $\text{mode}<\text{median}<\text{mean}$.
\end{proposition}

\begin{proof}
For real $k$, completing the square in the normal integral gives
$\E[e^{kY}] = \exp(k\mu + k^2\sigma^2/2)$. Writing $k\mu + k^2\sigma^2/2 = k(\mu+\sigma^2/2) + k(k-1)\sigma^2/2$,
the raw moments are
\begin{align}
  \E[X^k] &= m^k\,w^{k(k-1)/2}, \label{eq:lnraw}\\
  \text{so}\quad \E[X]&=m,\quad \E[X^2]=m^2w,\quad \E[X^3]=m^3w^3,\quad \E[X^4]=m^4w^6 . \nonumber
\end{align}
The central moments\ix{Central moments} $\mu_j=\E[(X-m)^j]$ follow from the binomial expansion:
\begin{align*}
  \mu_2 &= \E[X^2]-m^2 = m^2(w-1),\\
  \mu_3 &= \E[X^3]-3m\E[X^2]+2m^3 = m^3\bigl(w^3-3w+2\bigr) = m^3(w-1)^2(w+2),\\
  \mu_4 &= \E[X^4]-4m\E[X^3]+6m^2\E[X^2]-3m^4 = m^4\bigl(w^6-4w^3+6w-3\bigr)\\
        &= m^4(w-1)^2\bigl(w^4+2w^3+3w^2-3\bigr),
\end{align*}
where the factorisations can be checked by expanding the right-hand sides. Hence
\begin{equation*}
  S = \frac{\mu_3}{\mu_2^{3/2}} = \frac{(w-1)^2(w+2)}{(w-1)^{3/2}} = (w+2)\sqrt{w-1},
  \qquad
  K_e = \frac{\mu_4}{\mu_2^{2}} - 3 = w^4+2w^3+3w^2-6 ,
\end{equation*}
and $m$ cancels in both, so the shape depends on $\sigma$ alone. The exponential is increasing, so
$\mathbb P(X\le e^\mu)=\mathbb P(Y\le\mu)=\tfrac12$ and the median is $e^\mu$. The density of $X$ is
$g(x) = \varphi\bigl((\ln x-\mu)/\sigma\bigr)/(\sigma x)$ for $x>0$, where
$\varphi(z)=(2\pi)^{-1/2}e^{-z^2/2}$ is the standard normal density, and
\begin{equation*}
  \frac{d}{dx}\ln g(x) = -\frac{1}{x} - \frac{\ln x-\mu}{\sigma^2 x} = -\frac{\ln x-(\mu-\sigma^2)}{\sigma^2 x},
\end{equation*}
which is positive for $\ln x<\mu-\sigma^2$ and negative above it, so the mode is $e^{\mu-\sigma^2}$. The
ordering follows from $\mu-\sigma^2<\mu<\mu+\sigma^2/2$.
\end{proof}

Since $w>1$, the skewness\ix{Skewness} is positive and the distribution is skewed to the right. Because log returns add over time, a log return with parameters
$(\mu,\sigma^2)$ per year has parameters $(\mu h, \sigma^2 h)$ over $h$ years under independence, and the
skew grows with the horizon (\cref{fig:lognormal}). With $\mu = 6\%$ and $\sigma = 18\%$, the mean,
median and mode of terminal wealth per unit invested are \lnMeanOne, \lnMedianOne{} and \lnModeOne{}
after one year, and \lnMeanTwenty, \lnMedianTwenty{} and \lnModeTwenty{} after twenty years, where the
skewness has risen from \lnSkewOne{} to \lnSkewTwenty. Over long horizons\ix{Horizon} the mean is a poor guide to
what a typical investor will experience.

\begin{figure}[t]
  \centering
  \includegraphics[width=0.86\textwidth]{ch01/fig06_lognormal}
  \caption{Lognormal densities of terminal wealth per unit invested for log return parameters
  $\mu = 6\%$ and $\sigma = 18\%$ a year, at horizons of 1, 5 and 20 years.}
  \label{fig:lognormal}
\end{figure}

\FloatBarrier
% =====================================================================================
\section{Skew, fat tails and what a standard deviation promises}
\label{sec:moments}

\subsection{Third and fourth moments}

For a random return $R$ with mean $\mu = \E[R]$ and variance $\sigma^2 = \E[(R-\mu)^2]$, the
\emph{skewness\ix{Skewness}} and \emph{kurtosis\ix{Kurtosis}} are the standardised third and fourth central moments\ix{Central moments},
\begin{equation}\label{eq:skewkurt}
  S = \frac{\E[(R-\mu)^3]}{\sigma^3},\qquad
  K = \frac{\E[(R-\mu)^4]}{\sigma^4},\qquad K_e = K - 3 .
\end{equation}
The normal distribution has $S=0$ and $K=3$, and the \emph{excess kurtosis\ix{Excess kurtosis}}\index{Excess kurtosis} $K_e$ measures the fourth
standardised moment against that benchmark.  The uniform distribution has $K_e<0$, with a flat top and no tails to speak of; the scaled
Student $t$ law with five degrees of freedom\ix{Degrees of freedom} used below has $K_e>0$ and heavier tails than the normal.
Positive excess kurtosis does not fix the centre: many such laws have a taller, narrower peak, but
the shape near the mean varies from one to another. Positive skewness usually goes with a long right tail: frequent small losses and occasional
large gains, the profile of a lottery ticket or a long option. Negative skewness usually means the
reverse, steady small gains and occasional large losses.

From a sample $R_1,\dots,R_n$ with mean $\bar R$ and central moments
$m_k = n^{-1}\sum_i(R_i-\bar R)^k$, the usual estimators\ix{Estimation} are $s^2 = n\,m_2/(n-1)$ for the variance,
$\hat S = m_3/m_2^{3/2}$ and $\hat K_e = m_4/m_2^2 - 3$. The divisor $n-1$ makes $s^2$ unbiased. Small
sample corrections to $\hat S$ and $\hat K_e$ exist, and statistical packages differ in which they report,
so a skewness from one package and a skewness from another may not match for short series. This book
uses the uncorrected ratios above, which are the ones that enter the Jarque-Bera test.

\Cref{fig:kurtosis} compares four distributions scaled to unit variance. The uniform has
$K_e = -1.2$, the normal $0$, the Laplace $3$, and the Student $t$ with five degrees of freedom $6/(5-4)=\exkTfive$.
The differences are hard to see near the centre and large in the tails. The probability of a move
beyond four standard deviations is about \tailNormFour{} in 100{,}000 under the normal and about
\tailTFour{} in 100{,}000 under the scaled $t_5$, more than fifty times as large.

\begin{figure}[t]
  \centering
  \includegraphics[width=\textwidth]{ch01/fig07_kurtosis}
  \caption{(a) Four densities with mean zero and unit variance. (b) The right tails on a log scale,
  where differences in kurtosis become visible.}
  \label{fig:kurtosis}
\end{figure}

\subsection{Writing puts}

Option positions are a direct source of non-normal returns. Consider a fund that each month sells a
one-month put struck 5\% below the index level\ix{Index level}, holding cash equal to the strike\ix{Strike} as collateral. The
premium, priced with an implied volatility of 18\% and a continuously compounded rate of 3\% using the
formula derived in Chapter~5, is \putPrem{} per 100 of index. The premium is received at the start of the month and added to the collateral, and both earn the 3\%
rate. With strike $K$, premium $p$ and index level $S_T$ at expiry, the monthly return\ix{Monthly returns} on capital
equal to the strike is
\begin{equation*}
  R = \frac{(K+p)\,e^{0.03/12} - \max(K-S_T,0)}{K} - 1 .
\end{equation*}
The index starts at $S_0=100$ and follows a geometric Brownian motion with expected return 7\% and
volatility 16\% a year, $S_T = S_0\exp\{(0.07-\tfrac12 0.16^2)/12 + 0.16\,Z/\sqrt{12}\}$, so options are
priced a little rich relative to realised volatility. Taking $Z$ as 200{,}000 standard normal draws
from the chapter's shared generator, the first draws it makes, the simulated months give a mean monthly return of \pwMean\%, a standard deviation of \pwSD\%, skewness\ix{Skewness} of \pwSkew{}
and excess kurtosis\ix{Kurtosis}\ix{Excess kurtosis} of \pwExk. The strategy loses money in only \pwProbLoss\% of months, but those losses
are large (\cref{fig:putwriting}). A sample of 240 months, a long track record for most funds, gives
skewness \pwSkewS{} and excess kurtosis \pwExkS. A manager who reports only the mean and standard
deviation of such a strategy hides most of its risk.

\begin{figure}[t]
  \centering
  \includegraphics[width=\textwidth]{ch01/fig08_putwriting}
  \caption{Monthly returns of a cash-secured strategy that writes 5\% out-of-the-money puts, simulated
  under the index model in the text; no historical returns are used.
  (a) Histogram of 200{,}000 months, on a log density scale, with the normal density of the same mean and standard deviation.
  (b) Normal Q-Q plot of the first 240 months.}
  \label{fig:putwriting}
\end{figure}

\subsection{What one standard deviation promises}

If returns are normal, about 68\% of outcomes lie within one standard deviation\ix{Standard deviation} of the mean, 95\% within
two and 99.7\% within three. Without any assumption about the shape, Chebyshev\ix{Chebyshev inequality}'s inequality still gives
$\mathbb P(|R-\mu|\ge k\sigma)\le 1/k^2$. \Cref{tab:cheb} sets the two beside a pair of fat-tailed laws
scaled to unit variance, the Laplace and the Student $t$ with five degrees of freedom. Beyond two or
three standard deviations both put more probability than the normal does, though far less than
Chebyshev allows. Beyond one standard deviation they put less: \lapOutOne\% and \tOutOne\% against
\normOutOne\%, because the extra kurtosis comes with a taller, narrower centre. Neither the normal
column nor the bound says where an unknown distribution's tail probabilities lie. They are
reference points, and the distance between them shows how much the shape assumption is doing.

\begin{table}[t]
  \centering
  \caption{Probability (\%) of an outcome at least $k$ standard deviations from the mean, for three
  laws with unit variance and for Chebyshev's bound.}
  \label{tab:cheb}
  \small
  \begin{tabular}{@{}lrrrr@{}}
    \toprule
    $k$ & Normal & Laplace & Student $t_5$ & Chebyshev bound\\
    \midrule
    1 & \normOutOne & \lapOutOne & \tOutOne & \chebOne\\
    2 & \normOutTwo & \lapOutTwo & \tOutTwo & \chebTwo\\
    3 & \normOutThree & \lapOutThree & \tOutThree & \chebThree\\
    \bottomrule
  \end{tabular}
\end{table}

\FloatBarrier
% =====================================================================================
\section{Serial dependence: smoothing and the square root of time}
\label{sec:autocorr}

\subsection{Measuring serial dependence}

The \emph{autocorrelation\ix{Autocorrelation}} at lag $k$ is the correlation between a return and the return $k$ periods
earlier, $\rho_k = \Cov(R_t,R_{t-k})/\Var(R_t)$ for a stationary series, estimated by
$\hat\rho_k = \sum_{t=k+1}^n(R_t-\bar R)(R_{t-k}-\bar R)/\sum_t(R_t-\bar R)^2$. Under the null of
independent, identically distributed\ix{Identically distributed} returns with finite variance\ix{Finite variance}, $\hat\rho_k$ is approximately
$\N(0,1/n)$, which gives the $\pm1.96/\sqrt n$ bands used in the figures; with changing volatility or
other nonlinear dependence\ix{Dependence} those bands can be too narrow. The Durbin-Watson\ix{Durbin-Watson statistic} statistic
\parencite{durbin1950} was designed for the residuals $e_t$ of a regression and is
\begin{equation}\label{eq:dw}
  \mathrm{DW} = \frac{\sum_{t=2}^n (e_t-e_{t-1})^2}{\sum_{t=1}^n e_t^2}
  = \frac{\sum_{t\ge2} e_t^2 + \sum_{t\ge2} e_{t-1}^2 - 2\sum_{t\ge2} e_te_{t-1}}{\sum_t e_t^2}
  \approx 2(1-\hat\rho_1),
\end{equation}
where the approximation drops end terms of order $1/n$. Values near 2 indicate no first-order
autocorrelation, values toward 0 positive autocorrelation and values toward 4 negative autocorrelation. This chapter applies it descriptively to demeaned returns;
its tabulated critical values\ix{Critical value} belong to the regression setting.

\subsection{Appraisals and smoothing}

Appraisal\ix{Appraisal}-based returns, such as those of private real estate\ix{Real estate} or private equity NAVs, adjust slowly
to new information. A simple and widely used model \parencite{geltner1993} treats the reported return\ix{Reported returns}
$R^*_t$ as a weighted average\ix{Weighted average} of the economic return $R_t$ and last period's reported return:
\begin{equation}\label{eq:smooth}
  R^*_t = (1-\phi)R_t + \phi R^*_{t-1},\qquad 0\le\phi<1 .
\end{equation}

\begin{proposition}\label{prop:smooth}
If the economic returns are independent with mean $\mu$ and variance $\sigma^2$, the stationary reported
returns of \cref{eq:smooth} have mean $\mu$, variance $\sigma^2(1-\phi)/(1+\phi)$ and
autocorrelation\ix{Autocorrelation} $\phi^k$ at lag $k$.
\end{proposition}
\begin{proof}
Iterating \cref{eq:smooth} gives $R^*_t = (1-\phi)\sum_{j\ge0}\phi^jR_{t-j}$, whose mean is
$(1-\phi)\mu\sum_j\phi^j = \mu$. By independence its variance is
$(1-\phi)^2\sigma^2\sum_j\phi^{2j} = (1-\phi)^2\sigma^2/(1-\phi^2) = \sigma^2(1-\phi)/(1+\phi)$.
\Cref{eq:smooth} also gives $\Cov(R^*_t,R^*_{t-k}) = \phi\Cov(R^*_{t-1},R^*_{t-k})$ for $k\ge1$, since
$R_t$ is independent of the past, so the autocovariance shrinks by a factor $\phi$ per lag.
\end{proof}

Smoothing\ix{Return smoothing} leaves the mean return unchanged and hides volatility: the economic risk is still there,
but it no longer shows in the reported numbers. With $\phi = 0.6$, reported
volatility\ix{Measured volatility} is $\sqrt{0.4/1.6} = \smTheoryRatio\%$ of the true value. A simulation of 240 months, with economic
returns drawn as $\N(0.8\%,\,(4\%)^2)$ from the chapter's shared generator, has sample economic volatility of \smTrueSD\% and reported volatility of \smObsSD\%, a ratio of \smRatioSample{}
against the theoretical \smRatioTheory{} (\cref{fig:smoothing}). Annualised naively with $\sqrt{12}$, the
reported series suggests \smAnnObs\% a year when the simulated economic series annualises to \smAnnTrue\%;
the model value is \smAnnModel\%, and the gap between those two is sampling error in one finite path,
separate from the bias that smoothing introduces. The estimated
first-order autocorrelation\ix{First-order autocorrelation} is \smRhoOne, the Durbin-Watson\ix{Durbin-Watson statistic} statistic is \smDW{} and $2(1-\hat\rho_1)$ is
\smDWapprox. Inverting \cref{eq:smooth} recovers an estimate of the economic return,
\begin{equation}\label{eq:unsmooth}
  \hat R_t = \frac{R^*_t - \phi R^*_{t-1}}{1-\phi},
\end{equation}
and here, inverting with the true $\phi=0.6$, the unsmoothed\ix{Unsmoothing} series has volatility \smUnSD\% and lag-one autocorrelation \smUnRho. In practice
$\phi$ is unknown and is set to the estimated $\hat\rho_1$. That estimate identifies $\phi$ only under
the smoothing model with serially independent economic returns; if the economic returns are
themselves autocorrelated, it does not. \textcite{getmansky2004} develop a richer
version for hedge funds\ix{Hedge fund} holding illiquid assets. The same unsmoothing returns in Chapter~8 for risk
measures\ix{Risk measure} and in Chapter~10 for real estate indices.

\begin{figure}[t]
  \centering
  \includegraphics[width=\textwidth]{ch01/fig11_smoothing}
  \caption{(a) Simulated economic returns and the reported returns produced by smoothing with
  $\phi = 0.6$ (first 120 of 240 months). (b) Sample autocorrelations with the theoretical $\phi^k$; the shaded band is $\pm1.96/\sqrt{n}$. Both
  panels show one simulated path; the theoretical autocorrelations and the fall in volatility follow
  from the smoothing model.}
  \label{fig:smoothing}
\end{figure}

\subsection{Scaling risk with the horizon}

If one-period log returns\ix{Log return}, or additive changes such as yield changes, are independent and identically
distributed\ix{Identically distributed} with variance\ix{Variance} $\sigma^2$, the sum of $q$ of them has variance $q\sigma^2$, so volatility
scales with $\sqrt q$. For simple returns\ix{Simple return}, which compound, the rule holds only approximately. This is why monthly volatility is
annualised by multiplying by $\sqrt{12}$ and why, in the continuous-time models of Chapter~3, the
random part of a price change over a short interval $\delta t$ has size $\sqrt{\delta t}$ while the drift\ix{Drift}
has size $\delta t$ \parencite{wilmott2006}. The simplest discrete model writes the log return, or additive change, over a step of length
$\delta t$ as $R = \mu\,\delta t + \sigma Z\sqrt{\delta t}$ with $Z\sim\N(0,1)$, so that
the mean scales with $\delta t$ and the standard deviation with $\sqrt{\delta t}$. Over short steps the random
term dominates, which is why drift is so much harder to estimate than volatility. Autocorrelation\ix{Autocorrelation} breaks
the square root rule.

\begin{proposition}\label{prop:vr}
For a stationary\ix{Stationarity} series with autocorrelations $\rho_k$, the variance of a $q$-period sum relative to $q$
times the one-period variance is
\begin{equation}\label{eq:vr}
  VR(q) = \frac{\Var\bigl(\sum_{t=1}^q R_t\bigr)}{q\Var(R_t)} = 1 + 2\sum_{k=1}^{q-1}\Bigl(1-\frac kq\Bigr)\rho_k .
\end{equation}
For an AR(1)\ix{Autoregressive model} series with $\rho_k = \rho^k$ and $|\rho|<1$, $VR(q)\to(1+\rho)/(1-\rho)$ as $q\to\infty$.
\end{proposition}
\begin{proof}
$\Var(\sum_t R_t) = \sum_{s,t}\Cov(R_s,R_t) = \sigma^2\bigl[q + 2\sum_{k=1}^{q-1}(q-k)\rho_k\bigr]$, since
there are $q-k$ pairs at lag $k$. Dividing by $q\sigma^2$ gives \cref{eq:vr}. For $\rho_k=\rho^k$ with $|\rho|<1$, the terms
$\sum_k \rho^k \to \rho/(1-\rho)$ and $q^{-1}\sum_k k\rho^k \to 0$, so $VR \to 1 + 2\rho/(1-\rho) = (1+\rho)/(1-\rho)$.
\end{proof}

\Cref{fig:vr} plots $VR(q)$ for several values of $\rho$. For a smoothed series with $\rho = 0.6$ the variance of the 12-month sum is \vrTwelveSix{} times the variance it would have with no serial
correlation,
so the annual volatility implied by this stationary AR(1) model is \vrTwelveSixSqrt{} times the naive
figure. As the horizon grows the ratio of variances approaches $(1+0.6)/(1-0.6)=\vrLimitSix$, a
volatility multiplier of two. Applied to data, $\rho$ is estimated, so the correction carries its own
sampling error. The variance ratio\ix{Variance ratio}\index{Variance ratio} is also the basis of a classic test of the random walk
hypothesis \parencite{lo1988}.

\begin{figure}[t]
  \centering
  \includegraphics[width=0.86\textwidth]{ch01/fig12_variance_ratio}
  \caption{Variance ratio $VR(q)$ of an AR(1) series. Positive autocorrelation makes long-horizon risk
  larger than the square root rule suggests; negative autocorrelation makes it smaller.}
  \label{fig:vr}
\end{figure}

\FloatBarrier
% =====================================================================================
\section{What a short record says about normality}
\label{sec:normality}

\subsection{Sources of non-normality}

Three mechanisms account for most of the non-normal returns in this book: payoffs that are nonlinear in the underlying, prices
that adjust in jumps, and volatility that changes over time. Which way the skew points and how fat
the tails are depends on the mechanism and its parameters. A written option, as in the put-writing
example, gives a long left tail; a bought option or an incentive fee seen from the manager's side
gives a long right one. A jump can be a market event or the moment a stale price is finally marked.
A mixture of calm and turbulent periods is fat-tailed even when each period on its own is normal. Smoothing\ix{Return smoothing} is a
separate matter. A smoothed Gaussian series is still Gaussian, so smoothing alone cannot make returns non-normal. It does bias the measured variance and break the independence that the
Jarque-Bera test assumes, so moments and tests computed from appraisal-based series can mislead.

\subsection{The Jarque-Bera test}

\textcite{jarque1987} combine sample skewness and excess kurtosis\ix{Kurtosis}\ix{Excess kurtosis} into one statistic\ix{Test statistic},
\begin{equation}\label{eq:jb}
  \mathrm{JB} = \frac{n}{6}\left(\hat S^2 + \frac{\hat K_e^2}{4}\right).
\end{equation}
The weights $1/6$ and $1/24$ are the reciprocals of the asymptotic variances\ix{Asymptotic variance} of $\sqrt n\,\hat S$ and
$\sqrt n\,\hat K_e$ when the observations are independent and identically normal, the null under which
the $\chi^2_2$ reference distribution applies.

\begin{proposition}\label{prop:jb}
Let $R_1,\dots,R_n$ be independent draws from $\N(\mu,\sigma^2)$. As $n\to\infty$,
\begin{equation}\label{eq:jbclt}
  \sqrt n\begin{pmatrix}\hat S\\ \hat K_e\end{pmatrix} \xrightarrow{\;d\;}
  \N\!\left(\begin{pmatrix}0\\0\end{pmatrix},\begin{pmatrix}6&0\\0&24\end{pmatrix}\right),
  \qquad\text{hence}\qquad \mathrm{JB}\xrightarrow{\;d\;}\chi^2_2 .
\end{equation}
\end{proposition}

\begin{proof}
$\hat S$ and $\hat K_e$ do not change when every $R_i$ is shifted or rescaled by a positive constant, so we
may work with $z_i=(R_i-\mu)/\sigma\sim\N(0,1)$. Write $\overline{z^k}=n^{-1}\sum_i z_i^k$. The standard
normal moments are $\E[z^{2j}]=(2j-1)(2j-3)\cdots 1$, that is $1,3,15,105$ for $j=1,\dots,4$, and the odd
moments are zero. By the central limit theorem\ix{Central limit theorem} $\bar z$, $\overline{z^2}-1$ and $\overline{z^3}$ are of order
$n^{-1/2}$ in probability, written $O_p(n^{-1/2})$. Expanding the central moments\ix{Central moments} of the sample,
\begin{align*}
  m_2 &= \overline{z^2}-\bar z^2 = 1 + (\overline{z^2}-1) + O_p(n^{-1}),\\
  m_3 &= \overline{z^3} - 3\bar z\,\overline{z^2} + 2\bar z^3 = \overline{z^3} - 3\bar z + O_p(n^{-1}),\\
  m_4 &= \overline{z^4} - 4\bar z\,\overline{z^3} + 6\bar z^2\,\overline{z^2} - 3\bar z^4 = \overline{z^4} + O_p(n^{-1}),
\end{align*}
using $\bar z\,(\overline{z^2}-1)=O_p(n^{-1})$ in the second line. Since $m_2\to 1$ in probability,
\begin{align*}
  \sqrt n\,\hat S &= \frac{\sqrt n\, m_3}{m_2^{3/2}} = \sqrt n\;\overline{z^3-3z} + o_p(1),\\
  \sqrt n\,\hat K_e &= \frac{\sqrt n\,(m_4-3m_2^2)}{m_2^2}
   = \sqrt n\bigl[(\overline{z^4}-3) - 6(\overline{z^2}-1)\bigr] + o_p(1)
   = \sqrt n\;\overline{z^4-6z^2+3} + o_p(1),
\end{align*}
where the second line uses $m_2^2 = 1 + 2(\overline{z^2}-1) + O_p(n^{-1})$. Both leading terms are scaled
sample means\ix{Sample mean} of the i.i.d.\ variables $u=z^3-3z$ and $v=z^4-6z^2+3$, which have mean zero and
\begin{align*}
  \E[u^2] &= \E[z^6]-6\E[z^4]+9\E[z^2] = 15-18+9 = 6,\\
  \E[v^2] &= \E[z^8]-12\E[z^6]+42\E[z^4]-36\E[z^2]+9 = 105-180+126-36+9 = 24,\\
  \E[uv]  &= 0,
\end{align*}
the last because $uv$ is an odd polynomial in $z$. The multivariate central limit theorem and Slutsky's
lemma give \cref{eq:jbclt}. The two limits are jointly normal and uncorrelated, hence independent, and
$\mathrm{JB} = (\sqrt n\hat S/\sqrt6)^2 + (\sqrt n\hat K_e/\sqrt{24})^2$ is a continuous function of them;
by the continuous mapping theorem it converges to the sum of two independent squared standard normals,
which is $\chi^2_2$.
\end{proof}

The $\chi^2_2$ distribution has survival function $\mathbb P(\chi^2_2>x)=e^{-x/2}$, so the 5\% critical value\ix{Critical value} is
$-2\ln 0.05 = \jbCrit$ and the null of normality is rejected when $\mathrm{JB} > \jbCrit$. The listing
below computes the statistic and its asymptotic $p$-value:
\Needspace*{15\baselineskip}
\lstinputlisting[language=Python]{code/ch01/snippets/jb_test.py}
With seed 101 the normal sample gives $\mathrm{JB}=\jbSnipNorm$ ($p=\jbSnipNormP$), so normality is
not rejected, and the $t_5$ sample gives $\mathrm{JB}=\jbSnipT$ ($p=\jbSnipTP$), so it is.

The asymptotic approximation is poor in small samples, and alternative investment track records\ix{Track record} are
usually small samples. \Cref{fig:jb} reports a Monte Carlo\ix{Monte Carlo simulation} study with 20{,}000 samples at each size.
Under normal data the test rejects only \jbSizeTwenty\% of the time at $n=20$ and \jbSizeHundred\% at
$n=100$, against a nominal 5\%: it is conservative. Against Student $t_5$ data, whose excess kurtosis is
six, it rejects normality in \jbPowTwenty\% of samples at $n=20$, \jbPowHundred\% at $n=100$ and
\jbPowTwoFifty\% at $n=250$. Five years of monthly data is $n=60$; a failure to reject normality from
such a record says little about whether returns are normal, since at $n=60$ the test detects $t_5$ data in only \jbPowSixty\% of
samples. \textcite{thadewald2007} compare the Jarque-Bera\ix{Jarque-Bera test}
test with its competitors in more detail. The put-writing sample of 240 months, by contrast, gives
$\mathrm{JB} = \pwJB$, far beyond any critical value.

\begin{figure}[t]
  \centering
  \includegraphics[width=0.86\textwidth]{ch01/fig13_jb_size_power}
  \caption{Rejection rates of the Jarque-Bera test at the nominal 5\% level, from 20{,}000 simulated samples
  per size. Size under normal data is below 5\% in small samples; power against $t_5$ data is below
  two thirds at $n=100$ and passes 90\% only near $n=250$.}
  \label{fig:jb}
\end{figure}

\FloatBarrier
% =====================================================================================
\section{Co-movement: correlation, beta and diversification}
\label{sec:dependence}

\subsection{Three measures of co-movement}

The \emph{covariance} of returns $R_i$ and $R_j$ is $\sigma_{ij} = \E[(R_i-\mu_i)(R_j-\mu_j)]$, estimated
by $s_{ij} = (n-1)^{-1}\sum_t (R_{i,t}-\bar R_i)(R_{j,t}-\bar R_j)$. The \emph{Pearson correlation\ix{Correlation}}\index{Pearson correlation}
$\rho_{ij} = \sigma_{ij}/(\sigma_i\sigma_j)$ rescales covariance to $[-1,1]$ and measures the strength of
a linear relation. The \emph{Spearman\ix{Rank correlation} rank correlation}\index{Spearman rank correlation} \parencite{spearman1904} is the Pearson
correlation of the ranks of the observations. It measures the strength of any monotone relation and is
far less sensitive to a few extreme observations. When there are no ties it has the closed form
$\rho_S = 1 - 6\sum_t d_t^2/\bigl(n(n^2-1)\bigr)$, where $d_t$ is the difference between the two ranks of
observation $t$. The \emph{beta\ix{Beta}} of asset $i$ against a market $m$ is
the slope of the least squares regression of $R_i$ on $R_m$,
\begin{equation}\label{eq:beta}
  \beta_i = \frac{\Cov(R_i,R_m)}{\Var(R_m)} = \rho_{im}\frac{\sigma_i}{\sigma_m} .
\end{equation}
The two forms of \cref{eq:beta} are equal by the definition of $\rho_{im}$. That the ratio is the
regression slope is a short calculation.

\begin{proposition}\label{prop:ols}
For observations $(R_{m,t},R_{i,t})$, $t=1,\dots,n$, with $s_m^2>0$, the pair $(a,b)$ that minimises
$Q(a,b)=\sum_t (R_{i,t}-a-bR_{m,t})^2$ is unique and equals
\begin{equation}\label{eq:olsbeta}
  \hat b = \frac{\sum_t (R_{m,t}-\bar R_m)(R_{i,t}-\bar R_i)}{\sum_t (R_{m,t}-\bar R_m)^2} = \frac{s_{im}}{s_m^2},
  \qquad \hat a = \bar R_i - \hat b\,\bar R_m .
\end{equation}
\end{proposition}

\begin{proof}
$Q$ is a quadratic in $(a,b)$. Setting its partial derivatives to zero gives the normal equations\ix{Normal equations}
\begin{align*}
  \frac{\partial Q}{\partial a} &= -2\sum_t (R_{i,t}-a-bR_{m,t}) = 0
  &&\Longrightarrow\quad a = \bar R_i - b\bar R_m,\\
  \frac{\partial Q}{\partial b} &= -2\sum_t R_{m,t}(R_{i,t}-a-bR_{m,t}) = 0 .
\end{align*}
Substituting the first into the second and using $\sum_t (R_{m,t}-\bar R_m)=0$,
\begin{equation*}
  \sum_t R_{m,t}\bigl[(R_{i,t}-\bar R_i) - b(R_{m,t}-\bar R_m)\bigr]
  = \sum_t (R_{m,t}-\bar R_m)\bigl[(R_{i,t}-\bar R_i) - b(R_{m,t}-\bar R_m)\bigr] = 0,
\end{equation*}
which solves to $\hat b$ in \cref{eq:olsbeta}; the factors $n-1$ in $s_{im}$ and $s_m^2$ cancel. The Hessian
of $Q$ is $2\begin{psmallmatrix} n & \sum_t R_{m,t}\\ \sum_t R_{m,t} & \sum_t R_{m,t}^2\end{psmallmatrix}$, whose
determinant is $4n\sum_t (R_{m,t}-\bar R_m)^2>0$, so $Q$ is strictly convex and the stationary point is the
unique minimum.
\end{proof}

Correlation is symmetric and bounded; beta is neither. The beta of $i$ on $m$ is not the beta of $m$
on $i$, and it scales with the ratio of the two volatilities. A fund can be highly correlated with the market
and still have a low beta if it is much less volatile, as the next example shows.

\begin{example}{Beta and correlation from six months of returns}\label{ex:beta}
A fund and a market index return, in percent, $(1.2,-0.4,2.1,0.8,-1.0,1.1)$ and
$(2.0,-1.5,3.0,0.5,-2.5,1.5)$ over six months. The means are \exMeanF{} and \exMeanM. The sample
covariance is \exCov{} and the sample variance of the market is \exVarM, so the beta is
$\exCov/\exVarM = \exBeta$. The standard deviations are \exSdF{} and \exSdM, giving a Pearson correlation
of \exRho. The fund ranks the six months in exactly the same order as the market, so the Spearman
correlation is \exSpear. The fund moves with the market almost perfectly, but by about half as much.
\end{example}

\Cref{fig:correlation} shows two clouds of 300 bivariate normal draws with correlations of $\pm 0.7$,
and 60 points on the curve $y=e^{2.2x}$ with $x$ drawn uniformly on $[0,3]$, a relation that is
perfectly monotone but convex. For the convex relation the Pearson coefficient is \convexPearson{}
while Spearman's is \convexSpearman. A payoff\ix{Payoff} that is nonlinear but monotone can have perfect rank
correlation without perfect Pearson correlation; more complicated option payoffs need not be
monotone at all.

\begin{figure}[t]
  \centering
  \includegraphics[width=\textwidth]{ch01/fig09_correlation}
  \caption{Correlation pictures: two bivariate normal samples, and a monotone convex relation that the Pearson coefficient describes only in part, since it measures
  linear association.}
  \label{fig:correlation}
\end{figure}

\subsection{How far diversification goes}

\begin{proposition}\label{prop:div}
An equally weighted portfolio of $N$ assets, each with volatility $\sigma$ and pairwise correlation\ix{Correlation}
$\rho$, where $-1/(N-1)\le\rho\le1$ so that the correlation matrix is valid, has volatility
$\sigma_p = \sigma\sqrt{\rho + (1-\rho)/N}$. If $\rho\ge0$ is held fixed as $N\to\infty$, the volatility
falls to $\sigma\sqrt\rho$.
\end{proposition}
\begin{proof}
$\sigma_p^2 = \sum_i\sum_j N^{-2}\sigma_{ij} = N^{-2}\bigl[N\sigma^2 + N(N-1)\rho\sigma^2\bigr]
= \sigma^2\bigl[\rho + (1-\rho)/N\bigr]$. This is non-negative only if $\rho\ge-1/(N-1)$, so a fixed
correlation that is valid for every $N$ must be non-negative, and then $\sigma_p^2\to\sigma^2\rho$.
\end{proof}

When the common correlation is non-negative and stays fixed as the portfolio grows, the floor
$\sigma\sqrt\rho$ is the risk that diversification\ix{Diversification} cannot remove (\cref{fig:diversification}).
With $\sigma = 20\%$ and $\rho = 0.2$, ten assets bring volatility down to \divTenTwo\% and the floor is
\divFloorTwo\%. The appeal of many alternatives rests on a low correlation with the rest of the portfolio, and
smoothing\ix{Return smoothing} makes that correlation look lower than it is. Take the smoothing model \cref{eq:smooth}
with true returns $R_t=\beta M_t+e_t$, where the market return $M_t$ and the noise $e_t$ are
independent of each other and over time and the smoothing parameter is fixed.  Since $R^*_t=(1-\phi)\sum_{j\ge0}\phi^jR_{t-j}$ and only
the $j=0$ term moves with $M_t$, $\Cov(R^*_t,M_t)=(1-\phi)\beta\Var M$, so the measured beta\ix{Beta} is
$(1-\phi)\beta$. Dividing by the standard deviations, with $\Var R^*=\Var R\,(1-\phi)/(1+\phi)$ from
\cref{prop:smooth}, the measured correlation is the true one times
$(1-\phi)\sqrt{(1+\phi)/(1-\phi)}=\sqrt{1-\phi^2}$. At $\phi=\smAlpha$ the factors are
\smBetaFactor{} for beta and \smCorrFactor{} for correlation, and a simulation in the chapter's tests reproduces both. Other appraisal\ix{Appraisal} processes attenuate beta and
correlation by amounts that depend on how the appraisals are made.

\begin{figure}[t]
  \centering
  \includegraphics[width=0.86\textwidth]{ch01/fig10_diversification}
  \caption{Volatility of an equally weighted portfolio of assets with 20\% volatility as the number of
  assets grows, for four levels of pairwise correlation. Dotted lines mark the limits $\sigma\sqrt\rho$.}
  \label{fig:diversification}
\end{figure}

\FloatBarrier
% =====================================================================================
\section{Case study: daily changes in the five year Government of Canada yield}
\label{sec:case}

This section applies the tools of the chapter to real data: the daily five year Government of Canada\ix{Government of Canada bond}
benchmark bond yield\ix{Bond yield} (series \texttt{BD.CDN.5YR.DQ.YLD}) from the Bank of Canada \parencite{boc_valet},
taken from a public snapshot of Canadian policy rates\ix{Overnight rate}\ix{Policy rate} and bond yields \parencite{khan2026policy} that
also carries the target for the overnight rate (series \texttt{V39079}). The yields are the Bank's
published values; everything else in this section is computed from them by \texttt{goc\_case.py}.
The snapshot runs to \gocFileEnd, and the code cuts it at 30 September 2026, so the sample runs from
\gocStart{} to \gocEnd, the last September date with a yield, with \gocN{} daily observations. The file is archived with the code as \texttt{data/daily\_policy\_bond.csv}; its
SHA-256 digest is
\begin{center}\small\texttt{\gocShaA\gocShaB}\end{center}
and a test checks it, so the numbers below can be reproduced from exactly this file. Changes are measured between consecutive trading days\ix{Trading days}, so a change across a weekend
or holiday counts as one observation. The yield ranged from \gocYmin\% in \gocYminDate{} to
\gocYmax\% in \gocYmaxDate{} (\cref{fig:gocseries}). The target overnight rate stood at \polNow\% at
the end of the sample, the value the Bank's Valet service also reports for \gocEnd.

\begin{figure}[t]
  \centering
  \includegraphics[width=\textwidth]{ch01/fig14_goc_series}
  \caption{Top: the five year Government of Canada benchmark yield and the Bank of Canada target for the
  overnight rate. The target series in the snapshot begins on \polStart{} at \polFirst\%, the day the
  Bank cut it from 0.50\%, so the line starts there. Bottom: daily changes in the five year yield, in
  basis points.}
  \label{fig:gocseries}
\end{figure}

\paragraph{Moments and normality}
The \gocDN{} daily changes have a mean of \gocDmean{} basis points\ix{Basis point} and a standard deviation\ix{Standard deviation} of \gocDsd{}
basis points. Skewness is \gocDskew, close to zero, but excess kurtosis\ix{Kurtosis}\ix{Excess kurtosis} is \gocDexk. The Jarque-Bera\ix{Jarque-Bera test}
statistic\ix{Test statistic} is \gocJB. Under the asymptotic $\chi^2_2$ calibration, which assumes independent,
identically distributed normal changes, that is a clear rejection; volatility clustering\ix{Conditional heteroskedasticity} or serial
dependence\ix{Dependence} would call for a different calibration of the test. The fat tails\ix{Fat tails} are visible in
\cref{fig:gocdist}: there are \gocBeyondFour{} days with moves beyond four standard deviations, where a
normal distribution with the sample mean and standard deviation would expect \gocExpectFour. The largest rise was \gocMaxUp{} basis points on
\gocMaxUpDate{} and the largest fall \gocMaxDn{} basis points on \gocMaxDnDate. Heavy tails at the daily frequency are one of the stylised facts that
\textcite{cont2001} documents across many markets.

\begin{figure}[t]
  \centering
  \includegraphics[width=\textwidth]{ch01/fig15_goc_distribution}
  \caption{Standardised daily changes in the five year yield. (a) Histogram with the standard normal
  density. (b) Normal Q-Q plot; points above the line on the right and below it on the left are the
  fat tails.}
  \label{fig:gocdist}
\end{figure}

\paragraph{From yields to returns}
A yield change is not a return, but it maps into one. For a hypothetical five year par bond\ix{Par bond} priced each day at that day's starting yield, the one day
total return is approximately
$y\,\delta t - D\,\Delta y + \tfrac12 C(\Delta y)^2$: carry at the starting yield plus the duration\ix{Duration}
and convexity\ix{Convexity} terms of Chapter~6, ignoring roll-down and coupon reinvestment. Here $y$ is the yield at
the start of the day, $\delta t$ the calendar time\ix{Calendar time} to the next observation in years (actual/365), $D$
the modified duration\ix{Duration!modified} and $C$ the convexity. Carry is interest accrued over calendar days, so a Friday to Monday step earns three days
of it. The series is a hypothetical exposure\ix{Exposure}, rebuilt each day as a new five year par bond priced from
the observed yield; its returns approximate those of that exposure and are not the realised returns
of a bond held continuously or of a traded strategy. Annualising
treats each observation as one trading day, multiplying the mean by 252 and the standard deviation by
$\sqrt{252}$. With a median modified duration of \gocDur{} years, the return series has an annualised mean of \gocRmean\% and an annualised volatility\ix{Annualised volatility} of
\gocRsd\%, with excess kurtosis \gocRexk{} and a Jarque-Bera statistic of \gocRJB. The tails of the
yield changes carry over almost unchanged to the bond returns\ix{Bond return}.

\paragraph{Autocorrelation and the square root rule}
The lag\ix{Lag}-one autocorrelation\ix{Autocorrelation} of daily changes is \gocRhoOne, inside the $\pm\gocBand$ band, and the
Durbin-Watson\ix{Durbin-Watson statistic} statistic is \gocDW. Of the first ten lags only the tenth falls slightly outside the band
(panel (b) of \cref{fig:gocscaling}), about what chance alone would produce in ten tests at the 5\%
level. Unlike appraisal-based series, these daily yield changes show little sign of first-order serial
dependence. Panel (a) of \cref{fig:gocscaling} checks the square
root rule directly by comparing the standard deviation of non-overlapping $h$-day changes with
$\sqrt h$ times the daily standard deviation. The ratio is \gocRatioWeek{} at one week, \gocRatioMonth{}
at one month (standard error\ix{Standard error} about \gocSeMonth) and \gocRatioQuarter{} at one quarter (standard error
about \gocSeQuarter), all within two standard errors of one. These standard errors use the approximation for the sample
standard deviation of independent normal block changes; the heavy tails of this series, and any
dependence between blocks, make that approximation unreliable, and without further work they do not
say in which direction it errs.
Annualising daily volatility with $\sqrt{252}$ works for this series although the daily distribution
is far from normal. The two assumptions are separate. By \cref{eq:vr}, square root scaling holds exactly at every
horizon only when the autocorrelation is zero at every lag, and normality plays no part in it.
Autocorrelations that sum to zero are not enough: with $\rho_1=0.1$, $\rho_2=-0.1$ and none beyond,
\cref{eq:vr} gives a variance ratio\ix{Variance ratio} of \vrExTwo{} at two days, \vrExThree{} at three and \vrExTen{} at
ten, and it reaches one only in the limit.

\begin{figure}[!htb]
  \centering
  \includegraphics[width=\textwidth]{ch01/fig16_goc_scaling}
  \caption{(a) Standard deviation of non-overlapping $h$-day yield changes divided by $\sqrt h$ times the
  daily standard deviation, with bars of two standard errors from the independent normal
  approximation $\text{ratio}/\sqrt{2(m-1)}$, $m$ being the number of non-overlapping blocks. They indicate
  the size of the sampling error and are not confidence intervals calibrated to these heavy-tailed data. (b) Autocorrelation of daily changes with the
  $\pm1.96/\sqrt n$ band.}
  \label{fig:gocscaling}
\end{figure}

\FloatBarrier
% =====================================================================================
\section{Summary}
\label{sec:summary}

\begin{itemize}
  \item Compound through time with log returns, and aggregate across assets with simple returns.
  \item An IRR depends on the size and timing of cash flows. It can have several solutions or none, and
        it cannot reliably rank projects of different sizes or durations. An interim IRR rests on the
        current NAV, which is an estimate, and computing an IRR assumes nothing about reinvestment.
  \item Fees, the pace of deployment, expensed costs and cautious valuation can make a drawdown
        fund's early interim returns weak. The J-curve that results is a common pattern, though not every fund
        shows it.
  \item Positions with no entry price need a notional principal and a stated collateral level before
        a return can be computed. Under the stated assumptions, holding a fraction $c$ as collateral
        scales the derivative's gain or loss on capital by $1/c$, while the collateral interest is
        earned on the actual balance.
  \item Under independent normal log returns with positive variance, cumulative gross returns become
        more right-skewed as the horizon grows.
  \item Skewness and kurtosis matter most for strategies with option-like payoffs or prices that jump
        when they are finally marked. Beyond one standard deviation a fat-tailed law can carry less probability than the
        normal; beyond two or three it carries more.
  \item Smoothing lowers measured volatility, beta and correlation and induces autocorrelation. The
        variance ratio shows how much long-horizon risk is understated, and unsmoothing recovers an
        estimate of the economic returns under the smoothing model's assumptions.
  \item The Jarque-Bera test has little power against fat tails in samples of a few dozen
        observations, and a failure to reject is not evidence that returns are normal.
\end{itemize}

\FloatBarrier
% =====================================================================================
\section{Exercises}
\label{sec:exercises}

\begin{enumerate}[label=\textbf{\thechapter.\arabic*}, leftmargin=*]
  \item\label{exr:1} A bank quotes 6\% compounded monthly. Find the effective annual rate and the
        equivalent continuously compounded rate.
  \item\label{exr:2} A fully collateralised commodity futures position has a notional of \$5 million.
        Over a quarter the futures price falls 3\% and the collateral earns 1\% for the quarter on the
        actual collateral balance, with no financing, trading or other costs. Find the simple and log
        returns. Repeat for 50\% collateral.
  \item\label{exr:3} Show that a project with cash flows $-1$, $+3$, $-2$ has IRRs of 0\% and 100\%, and
        find the range of discount rates over which its NPV is positive.
  \item\label{exr:4} Two projects cost 100 each. One pays 160 in year 2; the other pays 130 in year 1.
        Find each IRR and the crossover rate, and say which project each rule prefers at a 10\% cost
        of capital.
  \item\label{exr:5} An investor puts 50 into a fund, which rises 10\%. The investor withdraws 30, and
        the fund then falls 10\%. Find the time-weighted and money-weighted returns per period and explain
        why they differ.
  \item\label{exr:6} Prove that the variance of the smoothed series in \cref{eq:smooth} is
        $\sigma^2(1-\phi)/(1+\phi)$ without iterating the recursion, by taking variances of both sides
        under stationarity.
  \item\label{exr:7} Using \cref{eq:vr}, find the variance ratio at $q=2$ for a series with $\rho_1 = -0.3$
        and interpret it.
  \item\label{exr:8} A sample of 60 monthly returns has skewness $-0.6$ and excess kurtosis $1.2$. Compute
        the Jarque-Bera statistic and state whether normality is rejected at 5\%. Use \cref{fig:jb}
        to say how much weight the result deserves, and give one reason the asymptotic 5\% critical value may mislead here, given that it assumes a
        large sample of independent, identically distributed returns.
\end{enumerate}

\subsection*{Computing projects}
\label{sec1:projects}
These use the data and code in the repository. Each is solved by a function in
\texttt{code/ch01/projects.py}, which \texttt{test\_projects.py} runs, so an answer can be
checked with \texttt{python code/ch01/projects.py} from the root of the repository. The
shipped series end in September 2026, and those that end earlier are noted.

\begin{enumerate}[label=\textbf{P1.\arabic*}, leftmargin=*, ref=P1.\arabic*]
\item Take month-end yields on Canadian 90 day bills from \texttt{data/goc\_bills\_bonds.csv}. Treat each
quote as a simple rate on an actual/365 basis for 91 days and convert it to an effective annual rate and
to a continuously compounded rate. How large is the largest gap between the quote and its effective
equivalent, in basis points, and in what month did it occur? Then compare the arithmetic and geometric
annualised returns of rolling the bill each month, and explain why they are so close.
\item Compute quarterly log changes of the Z.1 commercial real estate index in
\texttt{data/real\_assets\_quarterly.csv}. Report the sample moments and the Jarque-Bera test with
\texttt{sample\_moments} and \texttt{jarque\_bera}, estimate the first autocorrelation with \texttt{acf}, and
unsmooth the series with \texttt{unsmooth} at that coefficient. By what factor does the standard deviation
rise, and how does the result change if you start in 1996, when the index changes method?
\end{enumerate}

\subsection*{Answers}
\begin{enumerate}[label=\textbf{\thechapter.\arabic*}, leftmargin=*]
  \item $(1.005)^{12}-1 = 6.17\%$; $12\ln 1.005 = 5.985\%$.
  \item Full: $-3\% + 1\% = -2\%$, log return $\ln 0.98 = -2.020\%$. Half: $-6\% + 1\% = -5\%$, log return
        $\ln 0.95 = -5.129\%$.
  \item With $x = 1/(1+r)$, $-1 + 3x - 2x^2 = -(2x-1)(x-1) = 0$ gives $x=1$ and $x=1/2$, so $r = 0$ and
        $r = 100\%$. The NPV is positive for $1/2 < x < 1$, that is for $0 < r < 100\%$.
  \item $\sqrt{1.6}-1 = 26.49\%$ for the two-year project and $30\%$ for the one-year project. The
        crossover solves $160/(1+r)^2 = 130/(1+r)$, so $r = 160/130 - 1 = 23.08\%$, and below it the
        two-year project has the higher NPV. At 10\% its NPV is $160/1.21-100 = 32.23$ against
        $130/1.1-100 = 18.18$, so NPV prefers the two-year project while IRR prefers the one-year
        project. The rules conflict at any cost of capital below 23.08\%; between 23.08\% and 26.49\%
        both projects are worth doing and the one-year project also has the higher NPV.
  \item The balance is 55 after the first period and 25 after the withdrawal, and it ends at 22.5. TWR:
        $1.1\times0.9-1 = -1\%$ over two periods, or $\sqrt{0.99}-1 = -0.50\%$ per period. MWR:
        $-50 + 30/(1+r) + 22.5/(1+r)^2 = 0$, so $50(1+r)^2 - 30(1+r) - 22.5 = 0$ and
        $1+r = (30+\sqrt{900+4500})/100 = 1.0348$, an MWR of $3.48\%$ per period. The investor took money out
        before the loss, so their own experience beat the fund's.
  \item Under stationarity $\Var(R^*_t) = \Var(R^*_{t-1}) = v$ and $R_t$ is independent of $R^*_{t-1}$, so
        $v = (1-\phi)^2\sigma^2 + \phi^2 v$, which gives $v = (1-\phi)^2\sigma^2/(1-\phi^2) = \sigma^2(1-\phi)/(1+\phi)$.
  \item $VR(2) = 1 + 2(1/2)(-0.3) = 0.7$. Two-period variance is 30\% below twice the one-period variance; negative
        autocorrelation (mean reversion) makes longer-horizon risk smaller.
  \item $\mathrm{JB} = 10(0.36 + 0.36) = 7.2 > 5.991$, so normality is rejected at 5\%. The $\chi^2_2$ critical value is only asymptotic, but with 60 independent normal observations the
        actual rejection rate is below 5\% (\cref{fig:jb}), so over-rejection does not explain the
        result. A non-rejection would not have shown normality, because power at this size is low. The
        test also assumes independent, identically distributed returns: volatility clustering or
        smoothing changes the null distribution, so $p=e^{-3.6}\approx0.027$ should not be read literally.
\end{enumerate}

\FloatBarrier
% =====================================================================================
\section{Notes, sources and credits}
\label{sec:credits}

\paragraph{Further reading} \textcite{brealey2020} on IRR and NPV, \textcite{campbell1997} on
return statistics, \textcite{lo2002} on the sampling error of performance ratios, which returns in
Chapter~8, and \textcite{hull2022} on day count and compounding conventions.

\paragraph{Open source projects}
The code uses NumPy \parencite{numpy}, SciPy \parencite{scipy}, pandas \parencite{pandas}, Matplotlib
\parencite{matplotlib} and numpy-financial \parencite{numpyfinancial}. Each reference gives the version
installed and pinned in \texttt{code/requirements-lock.txt}, and a test checks that they agree. Selected calculations were also checked with statsmodels
\parencite{statsmodels}, R \parencite{rcore} and GNU Octave \parencite{octave}. The full list, with
links, is in \texttt{CREDITS.md}.

\paragraph{Data}
The case study uses Bank of Canada series \texttt{BD.CDN.5YR.DQ.YLD} and \texttt{V39079} as compiled in the snapshot of \textcite{khan2026policy}: the file \texttt{data/daily\_policy\_bond.csv}
at commit \texttt{815cb7b}. A copy is in this book's \texttt{data/} folder. Its SHA-256 digest, printed in
\cref{sec:case} and listed in \texttt{data/SOURCES.md}, can be checked with any SHA-256 tool
against the file at that commit, whose download address is also in \texttt{data/SOURCES.md}; a test
compares the copy with both, the second when a network is available. The file runs to \gocFileEnd;
the code reads it unchanged and drops every date after 30 September 2026, so no later observation
enters any result, in Python, R or Octave. The
observations themselves are the Bank of Canada's, published through its Valet service; the snapshot
only collects them.

\end{chapterrefs}
